\documentclass[11pt]{article}

\usepackage[letterpaper,margin=0.82in,headheight=15pt]{geometry}
\usepackage[T1]{fontenc}
\usepackage{lmodern}
\usepackage{microtype}
\usepackage{amsmath,amssymb,amsthm,bm,mathtools}
\usepackage{stmaryrd}
\usepackage{booktabs,longtable,tabularx,array,multirow}
\usepackage{enumitem}
\usepackage{xcolor}
\usepackage{fancyhdr}
\usepackage{graphicx}
\usepackage{tikz-cd}
\usepackage{hyperref}
\usepackage[nameinlink,noabbrev]{cleveref}

\definecolor{navy}{HTML}{17324D}
\definecolor{teal}{HTML}{147D92}
\definecolor{orange}{HTML}{D47A27}
\definecolor{softblue}{HTML}{EEF4F7}
\definecolor{softgray}{HTML}{F3F5F6}
\definecolor{darkgray}{HTML}{46545E}
\definecolor{softorange}{HTML}{FFF4E8}
\definecolor{softgreen}{HTML}{EDF7F0}
\definecolor{softred}{HTML}{FCEEEE}

\hypersetup{
  colorlinks=true,
  linkcolor=navy,
  citecolor=teal,
  urlcolor=teal,
  pdftitle={Volume X: Light-Sea Fock Sectors, Chiral Symmetry, PCAC, and Microscopic Antiquark GTMDs},
  pdfauthor={DeuteronWigner project}
}

\pagestyle{fancy}
\fancyhf{}
\fancyhead[L]{\small\color{darkgray}DeuteronWigner formalism - Volume X}
\fancyhead[R]{\small\color{darkgray}Light sea, chiral closure, and antiquark GTMDs}
\fancyfoot[L]{\small\color{darkgray}31 July 2026}
\fancyfoot[R]{\small\color{darkgray}\thepage}
\renewcommand{\headrulewidth}{0.4pt}

\setlist[itemize]{leftmargin=1.5em,itemsep=2pt,topsep=3pt}
\setlist[enumerate]{leftmargin=1.7em,itemsep=2pt,topsep=3pt}
\setlength{\parindent}{0pt}
\setlength{\parskip}{5pt}
\setlength{\emergencystretch}{2em}
\renewcommand{\arraystretch}{1.16}
\allowdisplaybreaks

\newtheorem{definition}{Definition}[section]
\newtheorem{axiom}[definition]{Axiom}
\newtheorem{principle}[definition]{Design principle}
\newtheorem{requirement}[definition]{Requirement}
\newtheorem{proposition}[definition]{Proposition}
\newtheorem{theorem}[definition]{Theorem}
\newtheorem{lemma}[definition]{Lemma}
\newtheorem{corollary}[definition]{Corollary}
\newtheorem{remark}[definition]{Remark}
\newtheorem{decision}[definition]{Model decision}

\crefname{definition}{definition}{definitions}
\crefname{axiom}{axiom}{axioms}
\crefname{principle}{design principle}{design principles}
\crefname{requirement}{requirement}{requirements}
\crefname{proposition}{proposition}{propositions}
\crefname{theorem}{theorem}{theorems}
\crefname{lemma}{lemma}{lemmas}
\crefname{corollary}{corollary}{corollaries}
\crefname{remark}{remark}{remarks}
\crefname{decision}{model decision}{model decisions}

\newcommand{\kT}{\bm{k}_{T}}
\newcommand{\pT}{\bm{p}_{T}}
\newcommand{\qT}{\bm{q}_{T}}
\newcommand{\DT}{\bm{\Delta}_{T}}
\newcommand{\bD}{\bm{b}_{\Delta}}
\newcommand{\bM}{\bm{b}_{\mathrm{TMD}}}
\newcommand{\dd}{\mathrm{d}}
\newcommand{\Tr}{\operatorname{Tr}}
\newcommand{\tr}{\operatorname{tr}}
\newcommand{\diag}{\operatorname{diag}}
\newcommand{\rank}{\operatorname{rank}}
\newcommand{\Span}{\operatorname{span}}
\newcommand{\Ker}{\operatorname{Ker}}
\newcommand{\Ran}{\operatorname{Ran}}
\newcommand{\Cov}{\operatorname{Cov}}
\newcommand{\Var}{\operatorname{Var}}
\newcommand{\Id}{\operatorname{Id}}
\newcommand{\Hom}{\operatorname{Hom}}
\newcommand{\End}{\operatorname{End}}
\newcommand{\Hil}{\mathcal H}
\newcommand{\LF}{\ensuremath{\mathrm{LF}}}
\newcommand{\BLFQ}{\ensuremath{\mathrm{BLFQ}}}
\newcommand{\TMD}{\ensuremath{\mathrm{TMD}}}
\newcommand{\GTMD}{\ensuremath{\mathrm{GTMD}}}
\newcommand{\GPD}{\ensuremath{\mathrm{GPD}}}
\newcommand{\QCD}{\ensuremath{\mathrm{QCD}}}
\newcommand{\EFT}{\ensuremath{\mathrm{EFT}}}
\newcommand{\TTN}{\ensuremath{\mathrm{TTN}}}
\newcommand{\TN}{\ensuremath{\mathrm{TN}}}
\newcommand{\PCAC}{\ensuremath{\mathrm{PCAC}}}
\newcommand{\WTI}{\ensuremath{\mathrm{WTI}}}
\newcommand{\cO}{\mathcal O}
\newcommand{\cM}{\mathcal M}
\newcommand{\cR}{\mathcal R}
\newcommand{\cS}{\mathcal S}
\newcommand{\cT}{\mathcal T}
\newcommand{\cV}{\mathcal V}
\newcommand{\cW}{\mathcal W}
\newcommand{\cE}{\mathcal E}
\newcommand{\cG}{\mathcal G}
\newcommand{\cK}{\mathcal K}
\newcommand{\cQ}{\mathcal Q}
\newcommand{\cP}{\mathcal P}
\newcommand{\cC}{\mathcal C}
\newcommand{\cZ}{\mathcal Z}
\newcommand{\cD}{\mathcal D}
\newcommand{\cN}{\mathcal N}
\newcommand{\cU}{\mathcal U}
\newcommand{\cF}{\mathcal F}
\newcommand{\cA}{\mathcal A}
\newcommand{\cB}{\mathcal B}
\newcommand{\cL}{\mathcal L}
\newcommand{\cJ}{\mathcal J}
\newcommand{\cI}{\mathcal I}
\newcommand{\eps}{\varepsilon}
\newcommand{\abs}[1]{\left|#1\right|}
\newcommand{\norm}[1]{\left\lVert#1\right\rVert}
\newcommand{\inner}[2]{\left\langle #1\middle|#2\right\rangle}
\newcommand{\avg}[1]{\left\langle #1\right\rangle}
\newcommand{\phys}{\mathrm{phys}}
\newcommand{\eff}{\mathrm{eff}}
\newcommand{\bare}{\mathrm{bare}}
\newcommand{\ind}{\mathrm{ind}}
\newcommand{\ct}{\mathrm{ct}}
\newcommand{\reg}{\mathrm{reg}}
\newcommand{\trunc}{\mathrm{trunc}}
\newcommand{\ren}{\mathrm{ren}}
\newcommand{\num}{\mathrm{num}}
\newcommand{\UV}{\mathrm{UV}}
\newcommand{\IR}{\mathrm{IR}}
\newcommand{\COM}{\mathrm{CM}}
\newcommand{\Lz}{L_z}
\newcommand{\Jz}{J^z}
\newcommand{\PV}{\operatorname{PV}}
\newcommand{\Hthree}{\Hil_{qqq}}
\newcommand{\Hfour}{\Hil_{qqqg}}
\newcommand{\Hfiveu}{\Hil_{qqqu\bar u}}
\newcommand{\Hfived}{\Hil_{qqqd\bar d}}
\newcommand{\AssumptionBundle}{\texttt{AssumptionBundle}}
\newcommand{\PredictionPlan}{\texttt{PredictionPlan}}
\newcommand{\StateBundle}{\texttt{H3MicroscopicStateBundle}}
\newcommand{\PCACManifest}{\texttt{PCACClosureManifest}}
\newcommand{\SeaLedger}{\texttt{SeaFlavorLedger}}
\newcommand{\CommonParent}{\texttt{MicroscopicCommonParentBundle}}
\newcommand{\WilsonHandoff}{\texttt{AntiquarkWilsonHandoff}}
\newcommand{\RenTrajectory}{\texttt{H3RenormalizationTrajectory}}
\newcommand{\ProvComplex}{\texttt{Provenance2Complex}}
\newcommand{\Amp}{\mathbf{Amp}}
\newcommand{\Match}{\mathbf{Match}}
\newcommand{\Red}{\mathbf{Red}}
\newcommand{\Proc}{\mathbf{Proc}}
\newcommand{\Infer}{\mathbf{Infer}}

\newcolumntype{Y}{>{\raggedright\arraybackslash}X}
\newcolumntype{L}[1]{>{\raggedright\arraybackslash}p{#1}}

\title{\vspace{-1.0cm}
{\color{navy}\bfseries Volume X: Light-Sea Fock Sectors, Chiral Symmetry,\\[2mm]
PCAC, and Microscopic Antiquark GTMDs}\\[4mm]
\large Fully antisymmetric five-parton dynamics, positive-
\(x\) antiquarks, axial and pseudoscalar currents, explicit/induced sea
matching, and the first common microscopic \(q\)--\(\bar q\)--\(g\) parent}
\author{DeuteronWigner project}
\date{31 July 2026}

\begin{document}
\maketitle

\begin{center}
\begin{minipage}{0.94\textwidth}
\colorbox{softblue}{\parbox{0.96\textwidth}{
\textbf{Status and purpose.}
Volumes~0--IX established the typed geometric architecture, regulated
light-front state spaces, common GTMD overlap maps, dynamical Wilson-line
pilots, matched nuclear and QCD interfaces, shared inference, a concrete
RSA--BLFQ Hamiltonian route, an assumption-conditioned tensor-network
compiler, an interacting valence state, and the first coupled
\(qqq\oplus qqqg\) microscopic eigenstate.  C10/H3 is the active
implementation package for the next required transition: explicit
\(u\bar u\) and \(d\bar d\) sectors with complete five-parton color
multiplicity, exact four-quark antisymmetry, canonical gluon splitting,
controlled chiral pair dynamics, Hamiltonian-consistent axial and
pseudoscalar operators, and direct positive-\(x\) antiquark matrix elements.
The present volume is the normative specification for that transition.  Its
central correction is that the quark-level nonsinglet axial Ward identity and
a hadronic pion-pole PCAC representation are matched descriptions of the same
divergence.  They may not be added as independent physical terms.  Likewise,
an explicit \(qqqq\bar q\) sector is not automatically an explicit pion
sector: a pion interpretation requires an additional spectral and operator
matching statement.  The resulting H3 state remains a finite-resolution
Hamiltonian-EFT state, not a continuum-QCD nucleon or a physical TMD.
}}
\end{minipage}
\end{center}

\begin{abstract}
A microscopic antiquark distribution cannot be generated by assigning a sea
probability to a valence wave function or by copying the negative-\(x\)
continuation of a quark curve.  It requires an interacting state containing
explicit antiquark creation operators, a complete color-singlet basis, exact
fermionic statistics, Hamiltonian vertices that create and annihilate the
pair, and currents and bilocal operators transformed consistently with the
same truncated Hamiltonian.  We formulate the first such DeuteronWigner
state in the direct sum
\(\Hthree\oplus\Hfour\oplus\Hfiveu\oplus\Hfived\).  The five-parton color
space contains three independent singlet multiplicities, one cluster-like and
two non-cluster couplings.  The four quarks are represented by an exterior
Fock construction or an exactly equivalent signed permutation basis before
operator assembly.  Canonical \(g\leftrightarrow q\bar q\) vertices and a
controlled isovector pseudoscalar/derivative interaction define distinct,
assumption-conditioned mechanisms.  Sector-dependent masses, pair vertices,
chiral coefficients, currents, and counterterms flow over a common resolution
tower.

The nonsinglet axial current is tested first through a finite-basis axial
Ward--Takahashi identity.  A second, hadronic PCAC representation isolates an
explicit or induced pion-pole channel and a non-pole remainder; an equivalence
cell prevents double counting between these two descriptions.  We define
axial, induced-pseudoscalar, and pseudoscalar form factors, a
Goldberger--Treiman-like holdout, and term-resolved closure defects.  The
coupled state is solved by exact diagonalization, matrix-free Krylov methods,
and a fermionic symmetry-adapted tree tensor network whose four Fock branches
retain color outer multiplicity, pair flavor, antiquark representation,
helicity, orbital angular momentum, and resolution identity.  Positive-\(x\)
antiquark GTMDs are obtained by active-antiquark overlaps derived from the
normal-ordered bilocal operator.  Vector, axial, and tensor moments use their
operator-specific charge-conjugation signatures rather than one universal
quark-minus-antiquark rule.  Quark, antiquark, and gluon GTMD, TMD, GPD, PDF,
current, energy--momentum, Wigner, and OAM reductions remain views of one
microscopic member.  A Feshbach comparison relates explicit pair sectors to
induced sea operators plus transformed observables and a visible remainder.
The final package specifies assumption plans, provenance two-cells, software
objects, analytic benchmarks, negative tests, and acceptance conditions for
the first common microscopic \(q\)--\(\bar q\)--\(g\) parent, while keeping
all QCD matching, evolution, nuclear, process, and inference gates closed.
\end{abstract}

\tableofcontents

\section{Scientific problem, scope, and interfaces}
\label{sec:scope}

\subsection{The missing microscopic degree of freedom}

The H2 state contains an interacting valence sector and a dynamical one-gluon
sector,
\begin{equation}
 \Hil_{\mathrm{H2}}
 =
 \Hthree\oplus\Hfour.
\end{equation}
It can generate regulated quark and gluon matrix elements and can provide the
microscopic amplitudes needed by the validated one-gluon Wilson engine.  It
cannot generate a direct positive-\(x\) antiquark overlap because no state in
its Hilbert space contains an antiquark creation operator.  The first H3 state
must therefore be
\begin{equation}
 \Hil_{\mathrm{H3}}
 =
 \Hthree
 \oplus\Hfour
 \oplus\Hfiveu
 \oplus\Hfived,
 \label{eq:h3-space}
\end{equation}
with one normalized proton or neutron eigenstate
\begin{equation}
 |N,\Lambda;r,\mathfrak A\rangle
 =
 |\Psi_{3}^{\Lambda}\rangle
 +|\Psi_{4g}^{\Lambda}\rangle
 +|\Psi_{5u}^{\Lambda}\rangle
 +|\Psi_{5d}^{\Lambda}\rangle.
 \label{eq:h3-state}
\end{equation}
Here \(r\) is the full regulator and truncation identity and
\(\mathfrak A\) is an immutable assumption bundle.

\begin{requirement}[Direct antiquark support]
A central positive-\(x\) antiquark result shall be produced only from an
active antiquark slot in a retained Fock sector or from a named matched
operator obtained by eliminating such a sector.  A copied negative-\(x\)
quark curve, an independently normalized sea profile, or an unlabeled scalar
probability is not a microscopic antiquark parent.
\end{requirement}

\subsection{Inputs from earlier volumes}

Volume X consumes the following validated interfaces:
\begin{enumerate}
 \item the regulated physical-sector, Hamiltonian, current, and truncation
       contracts of Volumes~I and VII;
 \item the zero-skewness momentum fibers, recoil authority, operator fibers,
       overlap kernels, and reduction maps of Volume~II;
 \item the Wilson-path, cut, phase-budget, and link-reversal interfaces of
       Volume~III;
 \item the matched nuclear, QCD, evolution, and inference readiness gates of
       Volumes~IV--VI;
 \item the symmetry-adapted tensor-network and assumption compiler of
       Volume~VIII;
 \item the coupled \(qqq\oplus qqqg\) Hamiltonian, sector renormalization,
       Ward benchmark, gluon/OAM exports, and microscopic Wilson handoff of
       Volume~IX.
\end{enumerate}
No interface is upgraded merely because H3 introduces a sea sector.  In
particular,
\begin{equation}
 \texttt{MICROSCOPIC\_Q\_QBAR\_G\_COMMON\_PARENT\_VALIDATED}
 \not\Rightarrow
 \texttt{PHYSICAL\_TMD}.
\end{equation}

\subsection{Primary model decisions}

\begin{decision}[Explicit light-sea route]
The first central H3 branch retains explicit \(u\bar u\) and \(d\bar d\)
Fock sectors.  Canonical gluon splitting and a controlled low-resolution
chiral pair kernel are separately typed Hamiltonian mechanisms.  An induced
sea or pion-cloud description is a matched alternative, not an additive
correction.
\end{decision}

\begin{decision}[Primary axial identity]
The primary microscopic identity is the nonsinglet quark-level axial
Ward--Takahashi identity.  A hadronic PCAC/pion-pole representation is a
matched decomposition of the same divergence into pole and non-pole pieces.
The two descriptions cannot be summed without an explicit subtraction map.
\end{decision}

\begin{decision}[Positive-\(x\) representation]
Quarks and antiquarks are stored as distinct positive-\(x\) objects.  Their
relation to a signed-\(x\) quark convention is implemented by a typed
charge-conjugation adapter that transforms the operator, endpoints, Wilson
path, color representation, Fourier phase, and helicity conventions.
\end{decision}

\subsection{Explicit nonclaims}

H3 does not establish:
\begin{itemize}
 \item a continuum or regulator-independent nucleon;
 \item complete spontaneous chiral-symmetry breaking on the light front;
 \item an explicit asymptotic pion state unless a pion Fock sector is added;
 \item the singlet axial anomaly or full non-Abelian Slavnov--Taylor system;
 \item a soft-subtracted or UV-matched TMD;
 \item physical antiquark phenomenology, nuclear matching, process
       factorization, evolution, or statistical inference.
\end{itemize}

\section{H3 Fock space, flavor, color, and statistics}
\label{sec:fock}

\subsection{Versioned sector identity}

At resolution \(r\), the sector set is
\begin{equation}
 \cF_r^{\mathrm{H3}}
 =
 \{3,4g,5u,5d\},
\end{equation}
where the labels denote
\begin{equation}
 3=qqq,
 \qquad
 4g=qqqg,
 \qquad
 5u=qqqu\bar u,
 \qquad
 5d=qqqd\bar d.
\end{equation}
Every sector identifier contains
\begin{equation}
 \nu=
 \left(
 n_q,n_{\bar q},n_g;
 B,Q,I,I_3,\Jz;
 K,N_{\max},b,\lambda;
 \cR_{\UV},\cR_{\IR};
 \text{endpoint},\text{zero-mode},\text{boundary}
 \right).
\end{equation}
The pair flavor is a sector identity, not a late projection label.

The intrinsic momentum manifold is
\begin{equation}
 \cM_{\nu}
 =
 \left\{
 (x_i,\bm\kappa_{iT})
 \ \middle|\
 x_i>0,
 \sum_i x_i=1,
 \sum_i\bm\kappa_{iT}=0
 \right\}.
 \label{eq:momentum-manifold}
\end{equation}
The basis excludes a gluon zero mode under the first declared policy; any
future zero-mode treatment is a distinct assumption branch.

\subsection{Flavor and charge symmetry}

Let the light-flavor doublet be
\begin{equation}
 \psi=
 \begin{pmatrix}u\\d\end{pmatrix},
 \qquad
 t^a_I=\frac{\tau^a}{2}.
\end{equation}
In the exact charge-symmetric branch, proton and neutron states form an
isospin doublet generated by the same Hamiltonian parameters and correlated
member identity.  The sea obeys
\begin{equation}
 \bar u_n=\bar d_p,
 \qquad
 \bar d_n=\bar u_p,
 \label{eq:sea-isospin}
\end{equation}
and therefore
\begin{equation}
 \left(\bar d-\bar u\right)_n
 =-
 \left(\bar d-\bar u\right)_p.
\end{equation}
Charge-symmetry breaking may be added only through explicit mass, QED, or
matched interaction operators and must preserve the common member map.

\subsection{Five-parton color multiplicity}

The five-parton color space is
\begin{equation}
 \cV_{5}^{c}
 =3^{\otimes4}\otimes\bar3.
\end{equation}
The singlet multiplicity follows from
\begin{equation}
 3^{\otimes3}
 =1\oplus8_{\rho}\oplus8_{\lambda}\oplus10.
\end{equation}
Multiplication by the fourth quark yields one fundamental \(3\) from each of
\(1\otimes3\), \(8_{\rho}\otimes3\), and
\(8_{\lambda}\otimes3\), while \(10\otimes3\) contains no fundamental.
Each fundamental then couples with \(\bar3\) to a singlet.  Hence
\begin{equation}
 \dim\operatorname{Inv}_{SU(3)}
 \left(3^{\otimes4}\otimes\bar3\right)
 =3.
 \label{eq:five-singlet-count}
\end{equation}

\begin{proposition}[Complete five-parton color basis]
For each fixed non-color configuration, the physical color subspace has an
orthonormal basis
\begin{equation}
 \left\{
 C_{1},C_{\rho},C_{\lambda}
 \right\},
\end{equation}
where \(C_1\) contains the conventional color-singlet three-quark cluster
coupled to a singlet pair, while \(C_{\rho}\) and \(C_{\lambda}\) are the two
independent octet-derived couplings.  Any single baryon--meson cluster tensor
spans at most one channel and is not a complete five-parton color basis.
\end{proposition}

The total color generator is
\begin{equation}
 \cQ^A
 =
 \sum_{i=1}^{4}t_i^A
 +t_{\bar q}^A,
 \qquad
 t_{\bar q}^A=-\left(t^A\right)^T,
\end{equation}
and every basis tensor satisfies
\begin{equation}
 \cQ^A C_{\alpha}=0,
 \qquad
 \inner{C_{\alpha}}{C_{\beta}}=\delta_{\alpha\beta}.
\end{equation}
Supported recoupling matrices between color fusion trees must be unitary and
must retain the outer-multiplicity index \(\alpha\in\{1,\rho,\lambda\}\).

\begin{requirement}[Color completeness]
All three singlet multiplicities shall survive through the basis,
Hamiltonian, currents, tensor network, GTMD operators, Feshbach maps, and
provenance records.  A missing channel is a basis defect, not an uncertainty
band.
\end{requirement}

\subsection{Exact four-quark antisymmetry}

Let \(h_q\) be the complete one-quark mode space, including momentum, flavor,
color, helicity, and orbital labels.  The mathematically primary four-quark
space is
\begin{equation}
 \bigwedge{}^4 h_q,
\end{equation}
so that the five-parton fermionic sector is
\begin{equation}
 \Hil_{5f}^{(r)}
 =
 \left[
 \bigwedge{}^4 h_q^{(r)}
 \right]
 \otimes h_{\bar q}^{(r)}.
 \label{eq:exterior-space}
\end{equation}
An occupation-number implementation is equivalent and preferred
computationally.  If an explicit slot basis is used, the antisymmetrizer is
\begin{equation}
 \cA_4
 =\frac{1}{4!}
 \sum_{\pi\in S_4}
 \operatorname{sgn}(\pi)U(\pi),
 \qquad
 \cA_4^2=\cA_4=\cA_4^{\dagger}.
 \label{eq:antisymmetrizer}
\end{equation}
Flavor is an internal label; exchanging a \(u\) and a \(d\) creation operator
still produces the fermionic sign while also exchanging their flavor labels.
A visually chosen baryon--meson partition does not weaken this requirement.

The tensor-network realization is \(\mathbb Z_2\)-graded.  For homogeneous
vectors,
\begin{equation}
 \tau(v\otimes w)
 =(-1)^{|v||w|}w\otimes v.
 \label{eq:fermionic-braiding}
\end{equation}
All leaf reorderings and recouplings carry an executable fermionic-swap
ledger.

\begin{requirement}[Statistics before dynamics]
Pauli-forbidden states shall be absent before Hamiltonian matrix elements are
assembled.  Post-diagonalization antisymmetrization, cluster-only statistics,
or sign repair at the observable layer is prohibited.
\end{requirement}

\subsection{Physical H3 projector}

The physical finite sector is
\begin{equation}
 \Hil_{\nu,\phys}^{(r)}
 =
 \cP_{\mathrm{singlet}}
 \cP_{\mathrm{fermion}}
 \cP_{B,Q,I,I_3,\Jz}
 \cP_{\COM}
 \Hil_{\nu,\mathrm{kin}}^{(r)}.
\end{equation}
Projector compatibility is tested by
\begin{equation}
 \delta_{AB}^{(r)}
 =\norm{[\cP_A,\cP_B]}.
\end{equation}
An exact symmetry pair must give zero.  A nonzero value is either a basis bug
or a named truncation defect.

\section{Coupled H3 Hamiltonian and dynamical mechanisms}
\label{sec:hamiltonian}

\subsection{Block mass operator}

The H3 mass-squared operator is
\begin{equation}
 \cM_{\mathrm{H3},r}^{2}
 =
 \begin{pmatrix}
  M_{33} & V_{3\leftarrow4g} & V_{3\leftarrow5u} & V_{3\leftarrow5d}\\
  V_{4g\leftarrow3} & M_{4g4g} & V_{4g\leftarrow5u} & V_{4g\leftarrow5d}\\
  V_{5u\leftarrow3} & V_{5u\leftarrow4g} & M_{5u5u} & V_{5u\leftarrow5d}\\
  V_{5d\leftarrow3} & V_{5d\leftarrow4g} & V_{5d\leftarrow5u} & M_{5d5d}
 \end{pmatrix}_{r},
 \label{eq:h3-hamiltonian}
\end{equation}
with
\begin{equation}
 V_{\alpha\leftarrow\beta}
 =V_{\beta\leftarrow\alpha}^{\dagger}.
\end{equation}
A block is present only when generated by a retained canonical or effective
operator.  An unsupported block is typed as absent with a symmetry or
truncation reason; it is not a silent numerical zero.

The diagonal blocks contain
\begin{equation}
 M_{\alpha\alpha,r}
 =M_{0,\alpha,r}
 +V_{\mathrm{conf},\alpha,r}^{\ind}
 +V_{\mathrm{diag},\alpha,r}^{\mathrm{can}}
 +V_{\chi,\alpha,r}^{\eff}
 +\delta M_{\alpha,r}^{2}
 +\Delta_{\alpha,r}^{2,\trunc}.
\end{equation}
Every term carries source/target sector, regulator, symmetry, parameter owner,
Hermitian partner, power-counting or naturalness metadata, and provenance.

\subsection{Canonical gluon splitting}

The canonical pair block connects the one-gluon sector to the pair sectors,
\begin{equation}
 V_{5f\leftarrow4g}^{\mathrm{pair}}
 :\Hfour\longrightarrow\Hil_{5f},
 \qquad f=u,d.
\end{equation}
A reduced finite-basis matrix element has the structure
\begin{equation}
 \langle\alpha_{5f}|
 V_{5f\leftarrow4g}^{\mathrm{pair}}
 |\alpha_{4g}\rangle
 =g_{45f,r}^{\bare}
 \cC_{45f}
 \cS_{45f}
 \cL_{45f}^{(r)}
 \cT_{45f}^{(r)}
 \cP_{45f}^{\mathrm{ferm}}
 \cR_{45f}^{(r)},
 \label{eq:pair-vertex}
\end{equation}
where the factors retain color, helicity, longitudinal, transverse,
fermionic-ordering, and regulator information.  Flavor is diagonal at the
canonical QCD vertex, but the multiplicity of available emitting and spectator
configurations produces nontrivial proton/neutron correlations.

The absorption block is generated as the exact adjoint.  Pair creation cannot
be represented by inserting a scalar probability into the state.

\subsection{Controlled chiral pair interaction}

At a low resolution, H3 permits an effective interaction in the isovector
scalar/pseudoscalar chiral channel.  One useful chirally paired operator basis
is
\begin{equation}
 \cO_{\chi,S}
 =\left(\bar\psi\psi\right)^2,
 \qquad
 \cO_{\chi,P}
 =\left(\bar\psi i\gamma_5\tau^a\psi\right)^2,
 \label{eq:chiral-pair-basis}
\end{equation}
or a nonlocal derivative realization matched to an auxiliary pion field.
The H3 Hamiltonian uses only a declared, regulator-smeared finite-basis kernel
from this operator space,
\begin{equation}
 V_{\chi,r}^{\eff}
 =G_{\chi,r}
 \left[
 c_{S,r}\cO_{\chi,S}^{(r)}
 +c_{P,r}\cO_{\chi,P}^{(r)}
 +\cdots
 \right].
 \label{eq:chiral-effective}
\end{equation}
Its purpose is to test a controlled chiral pair mechanism, not to assert that
a finite contact interaction is the unique low-energy reduction of QCD.  The
operator basis is motivated by standard chiral four-fermion constructions and
by explicit light-front nucleon--pion Hamiltonian studies
\cite{NJL1961,WuZhang2004,Du2020,Cao2026}.

The sector-changing contraction
\begin{equation}
 \Hthree\longleftrightarrow\Hil_{5f}
\end{equation}
selects an isovector pseudoscalar or derivative pair channel with explicit
parity, isospin, \(\Jz\), regulator, and naturalness identity.  The corresponding
axial and pseudoscalar current terms are owned by the same operator family.

\begin{requirement}[Chiral ownership]
A chiral Hamiltonian coefficient shall propagate to every state, axial,
pseudoscalar, pion-pole, sea-flavor, GTMD, and current contraction to which its
operator contributes.  It cannot become an independent normalization for
\(\bar d-\bar u\), an antiquark TMD, or one axial form factor.
\end{requirement}

\subsection{An explicit pair is not yet an explicit pion}

A five-parton amplitude may have a large projection onto an isovector
pseudoscalar \(q\bar q\) channel, but this does not by itself establish an
asymptotic pion state.  A pion interpretation requires a pole or spectral
statement, a normalization convention, and matching of the composite
interpolating operator.  We distinguish:
\begin{equation}
 \texttt{EXPLICIT\_QBARQ\_PAIR},
 \qquad
 \texttt{EXPLICIT\_PION\_FOCK\_SECTOR},
 \qquad
 \texttt{INDUCED\_PION\_POLE\_OPERATOR}.
\end{equation}
Only the first and third are available in H3.  The second belongs to a later
state-space extension.

\subsection{Assumption-conditioned branches}

The compiler shall support at least:
\begin{longtable}{L{0.19\textwidth}L{0.73\textwidth}}
\toprule
Plan & Dynamical content\\
\midrule
\endhead
H3-PLAN-A & Explicit \(u\bar u\) and \(d\bar d\); canonical
\(g\leftrightarrow q\bar q\); explicit chiral pair kernel; refitted induced
confinement; H2 instantaneous partners; no independent induced sea term.\\
H3-PLAN-B & Explicit pair sectors and canonical gluon splitting; chiral pair
kernel disabled; refitted induced confinement.\\
H2-REFERENCE & Read-only coupled \(qqq\oplus qqqg\) state; no explicit sea
prediction.\\
H3-INDUCED-SEA-REFERENCE & Pair sectors eliminated or absent; induced sea and
chiral operators plus transformed observables and a declared remainder.\\
\bottomrule
\end{longtable}
These are alternative complete theories at the declared truncation.  They may
be compared, but their amplitudes cannot be added.

\section{Sector-dependent renormalization over the H3 tower}
\label{sec:renormalization}

\subsection{Renormalization datum}

At resolution \(r\), define
\begin{equation}
 \mathfrak R_r^{\mathrm{H3}}
 =
 \left(
 \begin{array}{c}
 \cR_r,
 \theta_{3,r},\theta_{4g,r},\theta_{5u,r},\theta_{5d,r},\\
 \delta g_{34,r},\delta g_{45u,r},\delta g_{45d,r},
 \delta g_{\chi,r},\\
 \{\cC_i\},\cS_r,\cZ_r,\Delta_r
 \end{array}
 \right).
 \label{eq:h3-ren-data}
\end{equation}
The tower contains at least three nested points in longitudinal resolution,
transverse basis size, and sector basis dimension.  Bare masses, vertex
coefficients, current counterterms, and induced interactions may flow.  The
physical calibration conditions are held fixed.

\subsection{Restricted calibration and frozen holdouts}

A minimal validation calibration may use:
\begin{equation}
 M_N^2=M_{\mathrm{ref}}^2,
 \qquad
 F_1^p(0)=1,
 \qquad
 F_1^n(0)=0,
\end{equation}
plus one pair-vertex condition, one chiral/naturalness condition, and one axial
identity point.  The following remain frozen holdouts:
\begin{itemize}
 \item a second pair-vertex kinematic point;
 \item \(g_A\) or an independent axial-current combination;
 \item a second PCAC momentum transfer;
 \item a pion-pole or Goldberger--Treiman-like diagnostic;
 \item \(\bar d-\bar u\) or a related sea-flavor moment;
 \item a sea momentum, helicity, or OAM observable;
 \item a nonzero-transfer vector current and a rotational diagnostic.
\end{itemize}
A failed holdout revises the interaction, operator, truncation, or discrepancy
model.  It does not receive a new TMD-specific parameter.

\subsection{Identifiability and null directions}

Let \(R_i(\theta_r)\) be the calibration residuals and
\begin{equation}
 J_{ia}^{(r)}
 =\frac{\partial R_i}{\partial\theta_a}.
\end{equation}
The singular values of \(J^{(r)}\), the Hessian spectrum, profile likelihoods,
and naturalness combinations must be reported.  A null direction is an
identifiability statement.  It cannot be hidden by adding another fitted
observable after the holdout split is frozen.

Sector-dependent bare masses are regulator data.  They are not different
physical masses for a quark in different Fock sectors.  The common physical
light-quark mass convention and any \(m_d-m_u\) breaking operator remain
explicit.

\begin{requirement}[Trajectory, not point fit]
A successful fit at one finite basis is not a renormalized model.  H3 shall
report the flow of all bare, counterterm, induced, current, and chiral
coefficients and the stability of matched observables along the truncation
tower.
\end{requirement}

\section{Axial, pseudoscalar, and pion-pole operator system}
\label{sec:axial}

\subsection{Nonsinglet axial operators}

For two light flavors, define the isovector quark operators
\begin{align}
 A^{\mu,a}(x)
 &=
 \bar\psi(x)\gamma^{\mu}\gamma_5\frac{\tau^a}{2}\psi(x),
 \\
 P^a(x)
 &=
 \bar\psi(x)i\gamma_5\frac{\tau^a}{2}\psi(x).
 \label{eq:axial-pseudoscalar-def}
\end{align}
At the continuum operator level, the nonsinglet axial identity is
\cite{Nambu1960,Chen2021,Alexandrou2024}
\begin{equation}
 \partial_{\mu}A^{\mu,a}
 =2m_q P^a,
 \label{eq:qcd-pcac}
\end{equation}
up to the declared isospin-breaking generalization.  There is no nonsinglet
axial anomaly.  At finite light-front truncation, both sides must be reduced
with the same Hamiltonian transformation and regulator.

The corresponding amputated axial Ward--Takahashi identity is
\begin{equation}
 q_{\mu}\Gamma_{5}^{\mu,a}(p',p)
 =
 S^{-1}(p')\gamma_5\frac{\tau^a}{2}
 +
 \gamma_5\frac{\tau^a}{2}S^{-1}(p)
 -2m_q\Gamma_5^a(p',p),
 \label{eq:axial-wti}
\end{equation}
with all propagators, vertices, counterterms, and constrained-field pieces
defined in the same H3 model space.

\begin{decision}[Primary finite-basis axial test]
The first H3 closure test is the projected finite-basis form of
\cref{eq:axial-wti}.  The result may be labeled
\texttt{FINITE\_BASIS\_AXIAL\_WTI\_BENCHMARKED}; it cannot be labeled a
continuum chiral proof.
\end{decision}

\subsection{Hamiltonian-consistent current decomposition}

The H3 axial and pseudoscalar operators are
\begin{align}
 A_{r}^{\mu,a}
 &={}
 A_{(1),r}^{\mu,a}
 +A_{\mathrm{pair},r}^{\mu,a}
 +A_{\chi,r}^{\mu,a}
 +A_{\mathrm{inst},r}^{\mu,a}
 +\delta A_{\ct,r}^{\mu,a},
 \\
 P_{r}^{a}
 &={}
 P_{(1),r}^{a}
 +P_{\mathrm{pair},r}^{a}
 +P_{\chi,r}^{a}
 +P_{\mathrm{inst},r}^{a}
 +\delta P_{\ct,r}^{a}.
 \label{eq:h3-axial-decomp}
\end{align}
The one-body, pair, chiral, instantaneous, and counterterm pieces are not
separately observable.  They form one operator whose decomposition depends on
the resolution and model-space transformation.

The term-resolved defect is
\begin{equation}
 \delta_{5,r}^{a}
 =
 q_{\mu}A_{r}^{\mu,a}
 -2m_{q,r}P_{r}^{a},
 \label{eq:axial-defect}
\end{equation}
with a ledger
\begin{equation}
 \delta_{5,r}^{a}
 =
 \delta_{1b}
 +\delta_{\mathrm{pair}}
 +\delta_{\chi}
 +\delta_{\mathrm{inst}}
 +\delta_{\ct}
 +\delta_{\reg}
 +\delta_{\mathrm{basis}}.
\end{equation}
Removing any contribution required by the selected plan must produce a signed
nonzero ablation residual.

\subsection{Hadronic PCAC and pion-pole representation}

At the hadronic level, it is useful to isolate a pion-pole interpolating
operator \(\Phi^a\) and a non-pole remainder.  The quark pseudoscalar density
is decomposed by a matching map,
\begin{equation}
 2m_q P^a
 =
 f_{\pi}m_{\pi}^{2}\Phi^a
 +R_{\mathrm{nonpole}}^a.
 \label{eq:pion-pole-match}
\end{equation}
Inserting \cref{eq:pion-pole-match} into \cref{eq:qcd-pcac} gives the
hadronic PCAC representation
\begin{equation}
 \partial_{\mu}A^{\mu,a}
 =
 f_{\pi}m_{\pi}^{2}\Phi^a
 +R_{\mathrm{nonpole}}^a.
 \label{eq:hadronic-pcac}
\end{equation}
If H3 contains no explicit pion Fock sector, \(\Phi^a\) is an induced
composite or pole operator with a typed matching record.  It is not an
additional independent term to be added on top of \(2m_qP^a\).

\begin{proposition}[No double counting between WTI and pion-pole PCAC]
Let \(\cZ_{P\to\pi}\) denote the matching map implementing
\cref{eq:pion-pole-match}.  The quark-level and hadronic descriptions are
related by a commuting square
\begin{equation}
\begin{tikzcd}[column sep=large,row sep=large]
 2m_q P^a \arrow[r,equal] \arrow[d,"\cZ_{P\to\pi}"']
 & q_{\mu}A^{\mu,a} \arrow[d,equal]\\
 f_{\pi}m_{\pi}^{2}\Phi^a+R_{\mathrm{nonpole}}^a
 \arrow[r,equal]
 & q_{\mu}A^{\mu,a}.
\end{tikzcd}
\end{equation}
Selecting both \(2m_qP^a\) and
\(f_{\pi}m_{\pi}^{2}\Phi^a+R_{\mathrm{nonpole}}^a\) as additive physical
sources counts the same divergence twice.
\end{proposition}

The software shall therefore require a
\texttt{PCACRepresentation} field with mutually exclusive values such as
\begin{verbatim}
QUARK_LEVEL_AXIAL_WTI
HADRONIC_PION_POLE_DECOMPOSITION
\end{verbatim}
and an explicit matching two-cell between them.

\subsection{Nucleon axial and pseudoscalar form factors}

With \(q=p'-p\), \(Q^2=-q^2\ge0\), and the declared spinor convention,
define
\begin{align}
 \langle N(p')|A^{\mu,a}|N(p)\rangle
 &={}
 \bar u(p')
 \left[
 \gamma^{\mu}\gamma_5 G_A(Q^2)
 +\frac{q^{\mu}}{2M_N}\gamma_5 G_P(Q^2)
 \right]
 \frac{\tau^a}{2}u(p),
 \\
 \langle N(p')|P^a|N(p)\rangle
 &={}
 \bar u(p')i\gamma_5
 G_5(Q^2)
 \frac{\tau^a}{2}u(p).
 \label{eq:axial-form-factors}
\end{align}
The operator identity implies, after consistent treatment of phases and the
factor of \(i\), the generalized relation
\begin{equation}
 2M_N G_A(Q^2)
 -\frac{Q^2}{2M_N}G_P(Q^2)
 =2m_q G_5(Q^2)
 +\delta_{\mathrm{PCAC}}(Q^2).
 \label{eq:formfactor-pcac}
\end{equation}
The exact sign of the pseudoscalar convention is stored in the form-factor
adapter rather than inferred from a name.

The axial charge is
\begin{equation}
 g_A=G_A(0).
\end{equation}
A Goldberger--Treiman-like diagnostic is
\cite{GoldbergerTreiman1958}
\begin{equation}
 \Delta_{\mathrm{GT}}
 =1-
 \frac{g_A M_N}{f_{\pi}g_{\pi NN}},
 \label{eq:gt-discrepancy}
\end{equation}
where \(g_{\pi NN}\) is obtained from the declared pion-pole or interpolating
operator.  H3 may use this only as a holdout unless its ingredients were not
used in calibration.

\subsection{Singlet axial channel}

The flavor-singlet axial divergence contains the gluonic anomaly.  H3 does
not yet provide the complete anomaly operator and all required higher Fock
sectors.  Therefore
\begin{equation}
 \partial_{\mu}A_0^{\mu}
 =2m_qP_0+\frac{n_f\alpha_s}{4\pi}F^a_{\mu\nu}\widetilde F^{a\mu\nu}
\end{equation}
is recorded as a future operator-completeness requirement, not as an H3
closure identity.  The first PCAC benchmark is strictly nonsinglet.

\section{Exact, Krylov, and fermionic tensor-network solvers}
\label{sec:solvers}

\subsection{Solver hierarchy}

Every benchmark H3 block is solved by:
\begin{enumerate}
 \item dense Hermitian diagonalization;
 \item matrix-free Krylov application and eigensolution;
 \item exact full-support symmetry-adapted tensorization;
 \item a genuine variational TTN optimization in nested bond spaces.
\end{enumerate}
The first three are mutual oracles on small spaces.  The fourth is the
scalable approximation strategy.

The eigenproblem is
\begin{equation}
 \cM_{\mathrm{H3},r}^{2}
 |\Psi_{N,n}^{(r)}\rangle
 =M_{N,n}^{2(r)}|\Psi_{N,n}^{(r)}\rangle.
\end{equation}
The state is tracked across resolutions by exact quantum numbers, comparison
maps, principal angles, currents, sea-flavor fingerprints, and a deterministic
phase convention.

\subsection{Four-branch state tensor network}

The state network is a direct sum of symmetry-adapted fermionic tensor
networks \cite{Singh2010,Corboz2010},
\begin{equation}
 \cN_{\mathrm{state}}
 =
 \cN_{qqq}
 \oplus\cN_{qqqg}
 \oplus\cN_{qqqu\bar u}
 \oplus\cN_{qqqd\bar d}.
 \label{eq:h3-ttn}
\end{equation}
Every branch retains Fock identity at the root.  The five-parton branches
contain virtual labels for
\begin{equation}
 \left(
 R_c,\alpha_c;
 R_I,I_3;
 S,\Lz,\Jz;
 \pi_{S_4};
 f_{\bar q};
 r
 \right),
\end{equation}
where \(\alpha_c\in\{1,\rho,\lambda\}\) is the color outer multiplicity and
\(\pi_{S_4}\) is the four-quark permutation representation.

The tensor network is fermionic: its edge spaces are \(\mathbb Z_2\)-graded
and every recoupling or leaf ordering carries the braiding sign of
\cref{eq:fermionic-braiding}.  A bosonic tensor network with an after-the-fact
sign table is not an equivalent implementation unless the equivalence is
proved for every contraction path.

\subsection{Exact tensorization and variational convergence}

For each internal cut \(e\), the exact state has a block Schmidt
decomposition
\begin{equation}
 |\Psi\rangle
 =
 \sum_{\lambda,\alpha}
 s_{\lambda\alpha}^{(e)}
 |\lambda,\alpha\rangle_L
 \otimes
 |\bar\lambda,\alpha\rangle_R.
\end{equation}
The full allowed bond support must reproduce the exact state, currents,
PCAC matrix elements, flavor ledgers, and GTMD overlap blocks.

For nested variational spaces
\begin{equation}
 \cV_{\chi_1}\subseteq\cV_{\chi_2},
\end{equation}
the globally optimized energies satisfy
\begin{equation}
 E_{\chi_2}\le E_{\chi_1},
 \qquad
 E_{\chi}\ge E_{\mathrm{exact}}.
\end{equation}
Convergence must be reported in
\begin{equation}
 \left
 \{
 P_{4g},P_{5u},P_{5d},
 \avg{x_g},\avg{x_{\bar u}},\avg{x_{\bar d}},
 \Delta\bar u,\Delta\bar d,
 L_{\bar u},L_{\bar d},
 g_A,\delta_{\mathrm{PCAC}},
 W_{\bar q}
 \right
 \}.
\end{equation}
A low-rank state that reproduces the energy but erases sea flavor, antiquark
OAM, or PCAC cancellation is not converged.

\begin{requirement}[Observable-sensitive bond policy]
Bond selection shall reserve every symmetry block needed by the requested
observable before weight-based truncation is applied.  A unique antiquark,
color, OAM, or axial channel cannot be removed solely because its Schmidt
weight is small.
\end{requirement}

\section{Microscopic ledgers and light-sea structure}
\label{sec:ledgers}

\subsection{Probability and sector closure}

For every normalized state,
\begin{equation}
 P_3+P_{4g}+P_{5u}+P_{5d}=1,
 \label{eq:probability-ledger}
\end{equation}
where
\begin{equation}
 P_{\alpha}
 =\langle\Psi|\cP_{\alpha}|\Psi\rangle.
\end{equation}
These probabilities depend on regulator, resolution, and unitary
representation.  They are diagnostics, not universal observables.

\subsection{Valence flavor, charge, and baryon number}

For a proton member,
\begin{align}
 \int_0^1\dd x\,[u(x)-\bar u(x)]&=2,
 \\
 \int_0^1\dd x\,[d(x)-\bar d(x)]&=1.
 \label{eq:valence-ledger}
\end{align}
The charge and baryon ledgers are
\begin{align}
 Q_p
 &={}
 \frac23(N_u-N_{\bar u})
 -\frac13(N_d-N_{\bar d})=1,
 \\
 B_p
 &={}
 \frac13\sum_{q=u,d}(N_q-N_{\bar q})=1.
\end{align}
The neutron follows through the correlated isospin map, not a separately
fitted sea.

\subsection{Longitudinal momentum and canonical spin/OAM}

The momentum ledger is
\begin{equation}
 \avg{x_q}+\avg{x_{\bar q}}+\avg{x_g}=1.
 \label{eq:momentum-ledger}
\end{equation}
At the finite canonical light-front level,
\begin{equation}
 \frac12
 =
 \frac12\left(\Delta\Sigma_q+\Delta\Sigma_{\bar q}\right)
 +\Delta G
 +L_q+L_{\bar q}+L_g
 +\delta_{J,r}.
 \label{eq:spin-ledger}
\end{equation}
The decomposition is regulator, gauge, and truncation dependent.  Matching to
kinetic or gauge-invariant decompositions is deferred to the appropriate
operator scheme.

\subsection{Dynamical sea-flavor asymmetry}

Define
\begin{equation}
 A_{\bar q}(x)
 =\bar d_p(x)-\bar u_p(x),
 \qquad
 \mathcal A_{\bar q}
 =\int_0^1\dd x\,A_{\bar q}(x).
 \label{eq:sea-asymmetry}
\end{equation}
In H3-PLAN-A this asymmetry can emerge from the interplay of valence flavor,
canonical pair creation, and the isovector chiral kernel.  It receives no
independent normalization.  In a flavor-symmetric plan, a vanishing result is
reported as a consequence of that assumption bundle, not modified by hand.

Under exact charge symmetry,
\begin{equation}
 A_{\bar q}^{n}(x)=-A_{\bar q}^{p}(x).
\end{equation}
A useful collinear holdout is the Gottfried combination, under its standard
leading assumptions,
\begin{equation}
 S_G
 =\frac13+\frac23\int_0^1\dd x\,
 \left[\bar u_p(x)-\bar d_p(x)\right].
\end{equation}
H3 may compare with this relation only after the regulated distribution is
matched and evolved consistently; at the finite boundary it is a diagnostic
of flavor counting and shape.

\section{\texorpdfstring{Positive-\(x\)}{Positive-x} antiquark GTMD operators and overlaps}
\label{sec:antiquark}

\subsection{Operator origin from the field expansion}

The quark bilocal operator is
\begin{equation}
 \cO_q^{[\Gamma,\gamma]}(z)
 =
 \bar\psi_q(-z/2)\,
 \Gamma\,U_{\gamma}[-z/2,z/2]\,
 \psi_q(z/2).
 \label{eq:bilocal}
\end{equation}
Expanding the fields in creation and annihilation operators produces a quark
number term \(b^{\dagger}b\), an antiquark number term \(d^{\dagger}d\), and
pair-changing terms whose support depends on kinematics and operator order.
The positive-\(x\) antiquark kernel is derived from the normal-ordered
\(d^{\dagger}d\) contraction.  Its sign, spinor order, and endpoint order are
not guessed from the quark expression.

Define the typed active-antiquark spin kernel
\begin{equation}
 \overline{\cP}_{j}^{[\Gamma,\gamma]}
 =
 \operatorname{NormalOrder}_{d^{\dagger}d}
 \left[
 \cO_q^{[\Gamma,\gamma]}
 \right]_j,
 \label{eq:antiquark-projector}
\end{equation}
which is evaluated with anti-fundamental color transport and the declared
spinor convention.  This direct active-antiquark construction is the
microscopic counterpart of light-front meson-cloud and sea-spectator overlap
representations, but it retains one common Hamiltonian state rather than
fitting an independent sea wave function \cite{KoflerPasquini2017,Choudhary2023}.

\subsection{Zero-skewness active-antiquark overlap}

At \(\xi=0\), the regulated H3 antiquark GTMD is
\begin{align}
 W_{\bar q/N,\lambda'\lambda}^{[\Gamma,\gamma]}
 (x,\kT,\DT;r)
 ={}&
 \sum_{n\in\{5u,5d\}}
 \sum_{j\in\bar q}
 \int[\dd x]_n[\dd^2\bm\kappa_T]_n
 \nonumber\\
 &\times
 \delta(x-x_j)
 \delta^{(2)}(\kT-\bm\kappa_{jT})
 \psi_{n,\lambda'}^{*}
 \left(\{x_i,\bm\kappa_{iT}^{\mathrm{out}}\}\right)
 \overline{\cP}_{j}^{[\Gamma,\gamma]}
 \psi_{n,\lambda}
 \left(\{x_i,\bm\kappa_{iT}^{\mathrm{in}}\}\right).
 \label{eq:antiquark-overlap}
\end{align}
The initial and final intrinsic momenta use the single recoil authority from
Volume~II.  For active parton \(j\),
\begin{align}
 \bm\kappa_{jT}^{\mathrm{in}}
 &=\bm k_{jT}-\frac{1-x_j}{2}\DT,
 &
 \bm\kappa_{jT}^{\mathrm{out}}
 &=\bm k_{jT}+\frac{1-x_j}{2}\DT,
 \\
 \bm\kappa_{iT}^{\mathrm{in}}
 &=\bm k_{iT}+\frac{x_i}{2}\DT,
 &
 \bm\kappa_{iT}^{\mathrm{out}}
 &=\bm k_{iT}-\frac{x_i}{2}\DT
 \quad(i\ne j).
 \label{eq:antiquark-recoil}
\end{align}
No antiquark projector may rederive these shifts locally.

The basic H3 overlap is diagonal in retained parton number for the
zeroth-rescattering, \(\xi=0\), positive-\(x\) DGLAP region.  An off-diagonal
Fock contribution requires a named source such as pair-changing operator
content, ERBL kinematics, Wilson expansion, constrained fields, or a matched
induced operator.

\subsection{Charge-conjugation adapter}

Let \(C\) act on the bilinear by
\begin{equation}
 C\Gamma^T C^{-1}
 =\eta_C(\Gamma)\Gamma.
\end{equation}
A signed-\(x\) convention may be related to the positive-\(x\) antiquark object
only through an adapter that also transforms
\begin{equation}
 \left(
 \Gamma,
 \gamma,
 \text{endpoint order},
 3\leftrightarrow\bar3,
 \text{Fourier phase},
 \text{helicity phases}
 \right).
\end{equation}
For local leading operators, the charge-conjugation signatures imply the
schematic moment rules
\begin{align}
 \text{vector:}\quad
 &\int_0^1\dd x\,[q-\bar q],
 \\
 \text{axial:}\quad
 &\int_0^1\dd x\,[\Delta q+\Delta\bar q],
 \\
 \text{tensor:}\quad
 &\int_0^1\dd x\,[h_1^q-h_1^{\bar q}],
 \label{eq:c-signatures}
\end{align}
with exact normalization and phases stored in the operator registry.

\begin{requirement}[Operator-specific antiquark sign]
There is no universal quark-minus-antiquark rule.  Every local moment shall use
the charge-conjugation signature of its operator and a tested convention
adapter.
\end{requirement}

\subsection{Common quark--antiquark--gluon parent}

For one microscopic member \(s\), H3 exports
\begin{equation}
 \cW_s
 =
 \left(
 W_{q,s},W_{\bar q,s},W_{g,s}
 \right)
 \label{eq:common-parent-vector}
\end{equation}
with common Hamiltonian, state, regulator, recoil, member, and assumption
identities.  Separate species retain separate operator definitions; unification
occurs through the shared state and observable assembly.

A common-parent manifest records:
\begin{itemize}
 \item sector and active-slot ancestry;
 \item color outer multiplicity;
 \item spin, helicity, OAM, and target identity;
 \item operator, path, and matching status;
 \item all available TMD, GPD, PDF, current, EMT, Wigner, and OAM reductions;
 \item unavailable projections and the missing amplitude block.
\end{itemize}
A missing block returns \texttt{UNAVAILABLE\_AT\_H3\_TRUNCATION}, not a
physical zero.

\subsection{Forward Gram positivity}

At forward kinematics and at the unsubtracted amplitude level where a
probability interpretation applies,
\begin{equation}
 \Phi_{\alpha\beta}^{\bar q}
 =\sum_X A_{\alpha X}^{\bar q}
 \left(A_{\beta X}^{\bar q}\right)^*,
\end{equation}
so the full antiquark-target helicity matrix is positive semidefinite.  The
test is performed on the matrix parent.  It is not extended automatically to
off-forward GTMDs, Wigner quasidistributions, or finite-order subtracted
scheme objects.

\section{Matched reductions and local-operator closure}
\label{sec:reductions}

\subsection{Regulated reduction maps}

For any species \(a=q,\bar q,g\), define the regulated zero-skewness maps
\begin{align}
 \mathsf T_{\reg}[W_a]
 &=W_a(x,\kT,\DT=0),
 \\
 \mathsf G_{\reg}[W_a]
 &=\int\dd^2\kT\,W_a(x,\kT,\DT),
 \\
 \mathsf P_{\reg}[W_a]
 &=\int\dd^2\kT\,W_a(x,\kT,0),
 \\
 \mathsf F_{\reg}[W_a]
 &=\int\dd x\,\dd^2\kT\,W_a.
 \label{eq:regulated-reductions}
\end{align}
These produce a regulated TMD boundary, GPD boundary, PDF boundary, and local
matrix element.  They are not yet physical QCD objects in a specified
soft-subtracted scheme.

The same parent must satisfy the commuting closure
\begin{equation}
 \mathsf P_{\reg}[W_a]
 =
 \left.\mathsf G_{\reg}[W_a]\right|_{\DT=0}
 \label{eq:pdf-route-closure}
\end{equation}
within the numerical and regulator tolerance.  A species-specific correction
factor inserted only to enforce \cref{eq:pdf-route-closure} is prohibited.

\subsection{Vector, axial, tensor, and energy--momentum moments}

Let \(H^q\), \(\widetilde H^q\), and \(H_T^q\) denote the relevant regulated
quark GPD projections and similarly for positive-\(x\) antiquarks.  The local
operator signatures give
\begin{align}
 F_1^q(t)
 &=\int_0^1\dd x\,
 \left[
 H^q(x,0,t)-H^{\bar q}(x,0,t)
 \right],
 \label{eq:vector-moment}
 \\
 G_A^q(t)
 &=\int_0^1\dd x\,
 \left[
 \widetilde H^q(x,0,t)+\widetilde H^{\bar q}(x,0,t)
 \right],
 \label{eq:axial-moment}
 \\
 g_T^q(t)
 &=\int_0^1\dd x\,
 \left[
 H_T^q(x,0,t)-H_T^{\bar q}(x,0,t)
 \right].
 \label{eq:tensor-moment}
\end{align}
The equations are convention templates; the complete implementation includes
normalization, flavor generator, and spinor adapters.

The quark-plus-antiquark energy--momentum moment is
\begin{equation}
 A_{q+\bar q}(t)
 =\int_0^1\dd x\,x
 \left[
 H^q(x,0,t)+H^{\bar q}(x,0,t)
 \right],
\end{equation}
while the standard gluon convention may store
\begin{equation}
 H^g(x,0,0)=xg(x),
\end{equation}
so that
\begin{equation}
 A_g(t)=\int_0^1\dd x\,H^g(x,0,t).
\end{equation}
The registry must not copy the quark Mellin power into the gluon route.

\subsection{Axial and pseudoscalar closure from the common parent}

The axial GTMD/GPD projection must reduce to the same \(G_A\) used in the
finite-basis PCAC test.  Likewise the pseudoscalar operator and induced
pion-pole channel must be linked to the same H3 state and chiral parameters.
The required diagram is
\begin{equation}
\begin{tikzcd}[column sep=large,row sep=large]
 W_{q,\bar q}^{[\gamma^+\gamma_5]}
 \arrow[r,"\mathsf G_{\reg}"]
 \arrow[d,"\mathsf F_{\reg}"']
 & \widetilde H^{q,\bar q}_{\reg}
 \arrow[d,"\int\dd x\,\eta_C"]\\
 A_{\reg}^{+,a}
 \arrow[r,"\text{form-factor adapter}"']
 & G_A^a(t).
\end{tikzcd}
\label{eq:axial-commuting-square}
\end{equation}
A failure of this square identifies an inconsistent operator reduction,
current, state normalization, or matching adapter.

\subsection{Partonic Wigner and OAM reductions}

The partonic Wigner distribution is \cite{Meissner2009,Lorce2011}
\begin{equation}
 \rho_a(x,\kT,\bD)
 =
 \int\frac{\dd^2\DT}{(2\pi)^2}
 e^{-i\bD\cdot\DT}
 W_a(x,\kT,\DT),
 \label{eq:wigner}
\end{equation}
where \(\bD\) is conjugate to \(\DT\).  It is distinct from
\(\bM\), conjugate to \(\kT\), used in TMD matching and evolution.

A canonical OAM moment is
\begin{equation}
 L_z^a[\gamma]
 =
 \int\dd x\,\dd^2\kT\,\dd^2\bD
 \left(\bD\times\kT\right)_z
 \rho_{LU}^{a}(x,\kT,\bD;\gamma),
 \label{eq:oam-moment}
\end{equation}
with path, gauge, and scheme identity retained.  The antiquark OAM comes from
the active-antiquark parent and cannot be inferred by subtracting a quark OAM
curve from a total.

\subsection{Matching status}

Every H3 reduction carries at least
\begin{verbatim}
REGULATED_MICROSCOPIC_H3
LINK_SHORTENING_REQUIRED
UV_MATCHING_REQUIRED
RAPIDITY_SOFT_MATCHING_REQUIRED
NO_EVOLUTION_APPLIED
NO_PROCESS_MAP_APPLIED
\end{verbatim}
The small-\(\bM\) OPE, LF-to-QCD operator matching, soft subtraction, rapidity
renormalization, and process factorization are separate downstream maps.  A
successful regulated marginal does not satisfy them automatically.

\section{Explicit pair sectors versus induced sea operators}
\label{sec:feshbach}

\subsection{Feshbach reduction}

Let \(P\) project onto the retained \(qqq\oplus qqqg\) space and \(Q\) onto
the explicit pair sectors.  The exact finite-model effective Hamiltonian is
\begin{equation}
 H_{\eff}(E)
 =PHP+PHQ\frac{1}{E-QHQ}QHP.
 \label{eq:feshbach-h}
\end{equation}
The wave operator is
\begin{equation}
 \omega(E)
 =\frac{1}{E-QHQ}QHP.
\end{equation}
For any current or partonic operator, the model-space transformation must be
performed with the Hamiltonian \cite{Feshbach1958,Karmanov2008}.  Thus,
\begin{equation}
 O_{\eff}(E',E)
 =P\left[1+\omega^{\dagger}(E')\right]
 O
 \left[1+\omega(E)\right]P.
 \label{eq:feshbach-o}
\end{equation}
Using \(POP\) in place of \cref{eq:feshbach-o} is generally inconsistent.

The norm kernel is
\begin{equation}
 N_{\eff}(E)
 =P\left[1+\omega^{\dagger}(E)\omega(E)\right]P.
\end{equation}
Singularities and conditioning of \((E-QHQ)^{-1}\) are part of the
comparison manifest.

\subsection{Equivalence plus remainder}

At finite truncation, the explicit and induced descriptions satisfy
\begin{equation}
 \text{explicit pair theory}
 =
 \text{induced sea theory}
 +\text{transformed operators}
 +\text{remainder},
 \label{eq:explicit-induced}
\end{equation}
where the remainder records omitted energy dependence, basis support, higher
sectors, and matching defects.  A zero eigenvalue residual for selected states
does not justify setting the operator remainder to zero.

The provenance two-cell is represented schematically by
\begin{equation}
\begin{aligned}
 \Hil_{\mathrm{H3}}
 &\xrightarrow{\ \mathrm{solve\ and\ evaluate}\ }
 \{E,\langle O\rangle\}_{\mathrm{explicit}},\\
 \Big\downarrow\!\scriptstyle{\mathrm{Feshbach\ eliminate\ pair}}
 &\hspace{3.4cm}\Big\Updownarrow\!\scriptstyle{\delta_{\mathrm{rem}}}\\
 \Hil_{\mathrm{H2}}^{\eff}
 &\xrightarrow{\ H_{\eff},O_{\eff},N_{\eff}\ }
 \{E,\langle O\rangle\}_{\mathrm{induced}}.
\end{aligned}
\end{equation}
The compiler refuses a plan containing both explicit pair sectors and the
induced sea operator unless a separate overlap subtraction makes the union
well defined.

\subsection{Pion-cloud matching and overlap}

An induced pion-cloud model, an explicit \(q\bar q\) sector, and a future
explicit pion sector can share low-energy regions.  Their matched composition
must have the form
\begin{equation}
 W^{\mathrm{sea,matched}}
 =W^{\mathrm{explicit\ pair}}
 +W^{\mathrm{pion/induced}}
 -W^{\mathrm{overlap}}.
 \label{eq:sea-overlap}
\end{equation}
The overlap map is defined by a common asymptotic expansion or projection, not
by tuning the total normalization after addition.

\begin{requirement}[No sea double counting]
The composition engine shall reject explicit pair dynamics, an induced sea
operator, and an independent pion-cloud correction selected simultaneously
unless every shared region has an executable subtraction two-cell.
\end{requirement}

\section{Microscopic antiquark Wilson-line handoff}
\label{sec:wilson}

\subsection{Anti-fundamental transport}

For a quark Wilson line \(U_3[\gamma]\), the antiquark transport is
\begin{equation}
 U_{\bar3}[\gamma]
 =U_3[\gamma]^*,
\end{equation}
with generators
\begin{equation}
 t_{\bar3}^a=-\left(t_3^a\right)^T.
\end{equation}
The H3 Wilson handoff retains active antiquark flavor, helicity, pair-sector
ancestry, color outer multiplicity, OAM block, Hamiltonian coupling,
regulator, endpoint policy, and microscopic member.

\subsection{Charge conjugation and link reversal}

The antiunitary link-reversal adapter transforms the complete operator,
\begin{equation}
 \Theta:
 \left(
 P,P',\kT,\DT,\Gamma,\gamma,\bar3
 \right)
 \longmapsto
 \left(
 \bar P,\bar P',-\kT,-\DT,
 \Gamma_{\Theta},\gamma_{\Theta},\bar3
 \right),
\end{equation}
including complex conjugation and helicity phases.  A future/past sign cannot
be assigned from the process name alone.

\subsection{Cut support and readiness}

The H3 adapter reuses the validated path, pole, light-front resolvent,
cut-ledger, phase-budget, and soft-overlap types.  A discrete off-shell bound
state gives
\begin{equation}
 \operatorname{Im}\cA=0
\end{equation}
unless a separately declared continuum or finite-volume spectral rule
provides physical cut support.  A finite numerical \(\eps\) remains a
convergence oracle and never enters the physical result identity.

The highest H3 status is
\begin{verbatim}
MICROSCOPIC_ANTIQUARK_WILSON_INPUT_INTERFACE_VALIDATED
\end{verbatim}
not \texttt{WILSON\_READY}.  Physical naive-T-odd antiquark predictions still
require cut support, complete soft/rapidity accounting, LF-to-QCD matching,
and a qualified process map.

\section{Assumption compiler, tensor topology, and provenance}
\label{sec:compiler}

\subsection{H3 assumption bundle}

An H3 assumption bundle contains
\begin{equation}
 \mathfrak A_{\mathrm{H3}}
 =
 \left(
 \begin{array}{l}
 \text{sector set and basis resolution},\\
 \text{canonical pair and chiral interactions},\\
 \text{confinement and instantaneous-term branches},\\
 \text{renormalization and current conventions},\\
 \text{PCAC representation},\\
 \text{TTN topology and bond policy},\\
 \text{explicit/induced sea choice},\\
 \text{Wilson, matching, nuclear, process, and inference gates}
 \end{array}
 \right).
\end{equation}
The compiler validates exclusions and dependency closure before constructing
a Hamiltonian or network.

The chain is
\begin{equation}
\begin{aligned}
 \mathfrak A_{\mathrm{H3}}
 &\longrightarrow \PredictionPlan
 \longrightarrow
 \{\cM_{\mathrm{H3}}^2,A^{\mu},P,\cO_{\GTMD}\}
 \\
 &\longrightarrow \StateBundle
 \longrightarrow \CommonParent.
\end{aligned}
 \label{eq:compile-chain}
\end{equation}
A missing physical sector is never inserted silently.  A failed compilation
returns a minimal structured certificate.

\subsection{Three networks remain distinct}

The H3 calculation uses:
\begin{equation}
 \cN_{\mathrm{state}},
 \qquad
 \cN_{\mathrm{operator}},
 \qquad
 \cK_{\mathrm{prov}}.
\end{equation}
The state network represents amplitudes; the operator network represents
Hamiltonian, current, GTMD, and Wilson insertions; the provenance complex
represents assumptions, derivations, alternatives, and subtractions.
A contraction path is not a Wilson path, and a provenance path is not a
physical amplitude.

\subsection{H3 provenance two-complex}

The complex contains:
\begin{itemize}
 \item 0-cells: H3 states, sectors, operators, schemes, plans, and results;
 \item 1-cells: pair creation, chiral coupling, current insertion, reduction,
       Feshbach elimination, matching, and Wilson handoff;
 \item 2-cells: explicit/induced equivalence, WTI/PCAC equivalence,
       pion/sea overlap subtraction, current completion, and count-once cuts.
\end{itemize}
A nontrivial cycle in
\begin{equation}
 H_k(\cK_{\mathrm{prov}})
\end{equation}
is an audit signal for a missing relation, unresolved overlap, or genuinely
alternative mechanism.  It is not the value of a TMD or amplitude.

\begin{requirement}[Normalized composition]
Every authoritative H3 result shall have a normalized provenance path in
which equivalent descriptions are counted once, mutually exclusive plans do
not coexist, and every nonzero remainder is visible.
\end{requirement}

\section{Validation ladder and falsifiability}
\label{sec:validation}

The H3 validation ladder is ordered by dynamical depth.

\subsection{Algebraic and state gates}

The first gates are:
\begin{enumerate}
 \item five-parton singlet multiplicity three, generator annihilation,
       orthonormality, and recoupling unitarity;
 \item exact four-quark antisymmetry and Pauli exclusion;
 \item support, charge, baryon, isospin, \(\Jz\), and center-of-mass closure;
 \item block-Hamiltonian Hermiticity and matrix-free equality;
 \item exact/Krylov/full-bond TTN agreement;
 \item state tracking across resolution and avoided crossings.
\end{enumerate}

\subsection{Current and chiral gates}

The operator gates are:
\begin{enumerate}
 \item preservation of the H2 vector Ward benchmark;
 \item finite-basis nonsinglet axial WTI closure;
 \item equivalent pion-pole PCAC decomposition with no double counting;
 \item signed ablation residuals for every required current term;
 \item axial, pseudoscalar, and pion-pole form-factor consistency;
 \item a frozen second-point PCAC or Goldberger--Treiman-like holdout.
\end{enumerate}

\subsection{Common-parent gates}

The partonic gates are:
\begin{enumerate}
 \item direct positive-\(x\) \(\bar u\) and \(\bar d\) active-slot overlaps;
 \item common quark, antiquark, and gluon member identity;
 \item TMD/GPD/PDF/current route closure;
 \item valence, charge, baryon, momentum, helicity, and OAM ledgers;
 \item forward helicity-matrix positivity where applicable;
 \item explicit/induced sea equivalence plus a visible remainder;
 \item zero antiquark absorption without declared cut support.
\end{enumerate}

\subsection{Falsifiable holdouts}

A plan counts \(\bar d-\bar u\), \(g_A\), an axial radius, a pion-pole
quantity, or a sea OAM moment as a prediction only when its normalization and
shape were excluded from calibration.  Failure must revise the Hamiltonian,
chiral operator, current, basis, matching, or discrepancy model.  It cannot be
repaired by an antiquark-specific coefficient.

\section{Software realization and formal interfaces}
\label{sec:software}

\subsection{Core objects}

\begin{longtable}{L{0.29\textwidth}L{0.63\textwidth}}
\toprule
Object & Responsibility\\
\midrule
\endhead
\texttt{H3SectorSpace} & Versioned \(qqq\), \(qqqg\), \(qqqu\bar u\), and
\(qqqd\bar d\) spaces with physical projectors, regulator identity, support,
and exact quantum numbers.\\
\texttt{FivePartonColorBasis} & Three singlet multiplicities, total-generator
nullspace, phases, and recoupling matrices.\\
\texttt{FourQuarkAntisymmetry} & Exterior/occupation representation,
canonical creation order, signed permutations, and Pauli gates.\\
\texttt{PairCreationVertex} and \texttt{PairAnnihilationVertex} & Typed
\(g\leftrightarrow q\bar q\) blocks with color, flavor, helicity, momentum,
regulator, and generated adjoint.\\
\texttt{ChiralPairVertex} & Controlled isovector scalar/pseudoscalar or
derivative pair kernel with current ownership and naturalness metadata.\\
\texttt{H3Hamiltonian} & Hermitian four-sector mass operator, matrix-free
application, term registry, and discrepancy basis.\\
\texttt{H3Renormalization}\newline\texttt{Trajectory} & Resolution-indexed parameters,
conditions, null directions, covariance, flow, and holdouts.\\
\texttt{AxialCurrentOperator} & One-body, pair, chiral, instantaneous, and
counterterm axial-current blocks.\\
\texttt{PseudoscalarOperator} & Pseudoscalar density in the same model space
and convention.\\
\texttt{PionPoleOrInducedOperator} & Explicit or induced pion-pole channel,
non-pole remainder, and matching to the quark-level WTI.\\
\texttt{PCACClosureReport} & Quark-level WTI, hadronic PCAC decomposition,
term residuals, ablations, and holdouts.\\
\texttt{H3MicroscopicStateBundle} & Exact, Krylov, and TTN state identities,
sector amplitudes, currents, ledgers, readiness, and member correlation.\\
\texttt{SeaFlavorLedger} & \(u,\bar u,d,\bar d\) number, momentum, helicity,
OAM, charge, baryon, isospin, and asymmetry diagnostics.\\
\texttt{AntiquarkOverlap}\newline\texttt{Evaluator} & Positive-\(x\) active-antiquark overlap,
operator-derived spin kernel, recoil, and reduction routes.\\
\texttt{MicroscopicCommon}\newline\texttt{ParentBundle} & Shared \(q,\bar q,g\) parent,
operator identities, available reductions, and unavailable blocks.\\
\texttt{ExplicitInducedSea}\newline\texttt{Comparison} & Feshbach Hamiltonian, transformed
operators, norm kernel, remainder, conditioning, and provenance cell.\\
\texttt{AntiquarkWilson}\newline\texttt{Handoff} & Anti-fundamental path, active antiquark
identity, OAM/color ancestry, cut support, and phase budget.\\
\texttt{H3TensorNetwork}\newline\texttt{Manifest} & Fusion trees, fermionic swap ledger,
virtual representations, bond policy, solver history, and observable
convergence.\\
\bottomrule
\end{longtable}

\subsection{Fail-closed map classes}

The interfaces remain separated:
\begin{longtable}{L{0.15\textwidth}L{0.27\textwidth}L{0.48\textwidth}}
\toprule
Layer & H3 examples & Composition rule\\
\midrule
\endhead
\(\Amp\) & Pair/chiral vertices, Hamiltonian, state, axial current, GTMD
insertions & Linear maps compose only with compatible physical sectors and
operator identities.\\
\(\Match\) & Feshbach reduction, WTI-to-pion-pole adapter, LF-to-QCD future
map & Description-changing maps preserve declared matrix elements up to a
visible remainder.\\
\(\Red\) & Common GTMD reductions & Coordinate, rank,
operator, and charge-conjugation conventions must match.\\
\(\Proc\) & Future SIDIS, DY, DIS, and gluon-sensitive assembly & Requires
qualified factorization, hard/color maps, and matched TMDs.\\
\(\Infer\) & Future calibration and posterior prediction & Requires a complete
member bundle, covariance, likelihood, and frozen holdouts.\\
\bottomrule
\end{longtable}

An H3 object may be complete in \(\Amp\) and still fail every downstream
\(\Match\), \(\Proc\), and \(\Infer\) gate.

\subsection{Required property-based tests}

At minimum, the implementation shall test:
\begin{enumerate}
 \item five-parton singlet count, generator annihilation, orthonormality, and
       recoupling unitarity;
 \item exact antisymmetry and canonical creation-order signs;
 \item support, quantum-number, center-of-mass, and projector closure;
 \item Hamiltonian Hermiticity and generated adjoints for every pair block;
 \item pair and chiral selection rules;
 \item resolution-indexed renormalization conditions and holdout separation;
 \item vector Ward preservation and nonsinglet axial WTI closure;
 \item WTI/pion-pole PCAC equivalence without double counting;
 \item exact/Krylov/full-bond TTN agreement and variational monotonicity;
 \item observable-sensitive bond convergence;
 \item probability, valence, charge, baryon, momentum, and spin ledgers;
 \item direct positive-\(x\) antiquark active-slot selection;
 \item quark/antiquark/gluon same-member identity;
 \item vector, axial, tensor, and EMT charge-conjugation signatures;
 \item regulated TMD/GPD/PDF/current closure;
 \item explicit/induced Feshbach equivalence plus nonzero remainder;
 \item antiquark Wilson handoff and zero absorption without support;
 \item production, nuclear, evolution, process, and inference isolation.
\end{enumerate}

\section{Benchmark families}
\label{sec:benchmarks}

\begin{longtable}{L{0.12\textwidth}L{0.28\textwidth}L{0.50\textwidth}}
\toprule
ID & Benchmark & Required result\\
\midrule
\endhead
H3-A & Five-parton color completeness & Singlet multiplicity three,
generator annihilation, orthonormality, recoupling unitarity, and failure of a
cluster-only basis.\\
H3-B & Four-quark antisymmetry & Exact signed exchanges,
\(\cA_4^2=\cA_4=\cA_4^{\dagger}\), occupation-basis oracle, and rejection of
Pauli-forbidden states.\\
H3-C & Pair-vertex Hermiticity & \(qqqg\leftrightarrow qqqq\bar q\) blocks
for both flavors and all color multiplicities agree with their generated
adjoints on basis states and random complex superpositions.\\
H3-D & Chiral sector-changing kernel & Isospin, parity, and \(\Jz\) selection;
zero-coupling limit; parameter flow; explicit/induced branch separation.\\
H3-E & Sector renormalization & At least three resolutions, mass/charge
closure, pair/chiral flow, Jacobian and Hessian diagnostics, and frozen
holdouts.\\
H3-F & Vector Ward preservation & H2 Abelianized Ward benchmark remains
closed after pair-sector attachments and counterterms are added.\\
H3-G & Axial WTI and PCAC & Quark-level WTI, matched pion-pole
decomposition, term-resolved ablations, and a second-point holdout.\\
H3-H & Solver agreement & Dense, matrix-free Krylov, exact tensorization, and
full-bond TTN agree within declared tolerances.\\
H3-I & Observable-sensitive compression & A low bond loses sea flavor,
antiquark OAM/helicity, axial, or PCAC information while energy appears
comparatively accurate.\\
H3-J & Microscopic ledgers & Probability, valence flavor, charge, baryon,
momentum, canonical \(\Jz\), and proton/neutron partner relations close.\\
H3-K & Positive-\(x\) antiquark parent & Direct active-antiquark overlap,
TMD/GPD/PDF/current closure, no negative-\(x\) copy, and same-member
\(q,\bar q,g\) identity.\\
H3-L & Explicit/induced comparison & Hamiltonian and transformed-operator
equivalence, norm kernel, conditioning, visible remainder, and double-counting
rejection.\\
H3-M & Chiral sea asymmetry & Plan-dependent \(\bar d-\bar u\), no independent
normalization, correlated neutron relation, and axial/pion-pole holdouts.\\
H3-N & Antiquark Wilson handoff & Anti-fundamental transport, complete state
identity, charge-conjugation/link conventions, and zero absorption without
support.\\
H3-O & Assumption compiler & Deterministic PLAN-A, PLAN-B, H2 reference, and
induced-sea reference; incompatible combinations fail before assembly.\\
H3-P & WTI/PCAC representation cell & Quark-level and pion-pole routes agree
after matching; additive use yields a deliberate factor-of-two or equivalent
nonzero defect.\\
\bottomrule
\end{longtable}

\section{Stable formal requirements}
\label{sec:requirements}

The implementation shall use stable machine-readable requirement IDs.  The
minimum set is:
\begin{longtable}{L{0.23\textwidth}L{0.69\textwidth}}
\toprule
ID & Requirement\\
\midrule
\endhead
\texttt{V10.STATE.1} & H3 state contains distinct \(qqq\), \(qqqg\),
\(qqqu\bar u\), and \(qqqd\bar d\) sectors.\\
\texttt{V10.STATE.2} & Every sector carries support, quantum numbers,
regulator, boundary, member, and assumption identities.\\
\texttt{V10.COLOR.1} & Five-parton singlet multiplicity is exactly three.\\
\texttt{V10.COLOR.2} & All color outer multiplicities remain visible through
state and operator exports.\\
\texttt{V10.FERMION.1} & Four-quark antisymmetry is exact before Hamiltonian
assembly.\\
\texttt{V10.FERMION.2} & Tensor-network recoupling uses an executable
fermionic-swap ledger.\\
\texttt{V10.HAM.1} & H3 Hamiltonian is Hermitian with named absent blocks.\\
\texttt{V10.HAM.2} & Canonical \(g\leftrightarrow q\bar q\) blocks and
adjoints are generated from one identity.\\
\texttt{V10.CHIRAL.1} & Chiral pair interaction has explicit operator,
selection, regulator, and parameter ownership.\\
\texttt{V10.CHIRAL.2} & Explicit pair, explicit pion, and induced pion-pole
identities remain distinct.\\
\texttt{V10.PLAN.1} & H3 assumption branches are content addressed and
mutually exclusive where required.\\
\texttt{V10.REN.1} & Sector, pair, chiral, current, and counterterm parameters
flow over a multi-resolution tower.\\
\texttt{V10.REN.2} & Null directions and holdouts remain visible.\\
\texttt{V10.AXIAL.1} & Primary nonsinglet axial WTI is implemented with one
Hamiltonian/regulator identity.\\
\texttt{V10.AXIAL.2} & Hadronic pion-pole PCAC is a matched alternative
decomposition.\\
\texttt{V10.AXIAL.3} & Singlet anomaly remains explicitly unresolved.\\
\texttt{V10.CURRENT.1} & Vector Ward benchmark remains passing.\\
\texttt{V10.CURRENT.2} & Axial, pseudoscalar, pair, chiral, instantaneous, and
counterterm pieces share ownership.\\
\texttt{V10.TTN.1} & Exact/Krylov/full-bond TTN states agree.\\
\texttt{V10.TTN.2} & Variational bond convergence includes sea and axial
observables.\\
\texttt{V10.LEDGER.1} & Probability, valence, charge, baryon, momentum, and
canonical spin ledgers close.\\
\texttt{V10.LEDGER.2} & Proton and neutron are correlated microscopic
members.\\
\texttt{V10.SEA.1} & \(\bar u\) and \(\bar d\) arise from direct positive-
\(x\) active slots.\\
\texttt{V10.SEA.2} & Sea asymmetry is a shared-state prediction, not a fitted
normalization.\\
\texttt{V10.GTMD.1} & Common regulated \(q,\bar q,g\) parent uses one state
and recoil authority.\\
\texttt{V10.GTMD.2} & Missing amplitude content is unavailable, not zero.\\
\texttt{V10.RED.1} & TMD/GPD/PDF/current routes commute at the regulated
level.\\
\texttt{V10.RED.2} & Vector, axial, tensor, and EMT moments use operator-
specific charge-conjugation signatures.\\
\texttt{V10.WIGNER.1} & \(\bD\) and \(\bM\) remain distinct.\\
\texttt{V10.POS.1} & Positivity is tested only on applicable forward Gram
matrices.\\
\texttt{V10.FESH.1} & Explicit and induced sea descriptions include
transformed operators, norm kernel, and remainder.\\
\texttt{V10.PROV.1} & Explicit pair and induced replacement cannot be active
together without subtraction.\\
\texttt{V10.PROV.2} & WTI and pion-pole PCAC are joined by an equivalence
cell, not an additive edge.\\
\texttt{V10.WILSON.1} & Antiquark handoff uses anti-fundamental transport and
complete microscopic ancestry.\\
\texttt{V10.WILSON.2} & No absorption occurs without declared physical cut
support.\\
\texttt{V10.MATCH.1} & UV, rapidity/soft, link-shortening, and LF-to-QCD
matching remain explicit requirements.\\
\texttt{V10.GATE.1} & Nuclear, evolution, process, and inference consumers
fail closed.\\
\texttt{V10.REGRESS.1} & Accepted production artifacts and registries remain
byte-identical.\\
\texttt{V10.FALSIFY.1} & Failed holdouts revise shared physics rather than a
named antiquark/TMD coefficient.\\
\bottomrule
\end{longtable}

\section{Staged research and implementation program}
\label{sec:stages}

\subsection{H3.0: exact state-space extension}

Implement the two flavor-resolved pair sectors, complete color multiplicity,
exact antisymmetry, and physical projectors.  No dynamics are added until the
basis gates pass.

\subsection{H3.1: canonical pair dynamics}

Add \(g\leftrightarrow q\bar q\) vertices and adjoints, extend the H2
Hamiltonian, and validate block Hermiticity and matrix-free action.

\subsection{H3.2: controlled chiral and current sector}

Add the chiral pair kernel, axial and pseudoscalar operators, the primary
axial WTI, and the matched pion-pole PCAC representation.  Freeze holdouts
before final coefficient optimization.

\subsection{H3.3: coupled-state and TTN tower}

Generate multi-resolution exact, Krylov, and variational TTN states; track the
physical nucleon; report parameter and bond convergence in sea and axial
observables.

\subsection{H3.4: common microscopic partonic parent}

Evaluate positive-\(x\) quark, antiquark, and gluon GTMD helicity blocks,
regulated reductions, local moments, Wigner/OAM quantities, and all ledgers.
This stage supplies the first candidate state for the future microscopic
replacement of the analytic C3/C4 parents.

\subsection{H3.5: explicit/induced and Wilson interfaces}

Complete Feshbach comparisons, normalized provenance cells, antiquark Wilson
handoff, and all downstream readiness gates.  No production replacement is
made in Volume X.

\subsection{Next package after H3}

Once C10/H3 reports, the next computational package should be
\begin{quote}
\textbf{H4: microscopic nonzero-transfer quark, antiquark, and gluon GTMD
helicity matrices from the common H3 eigenstate; complete T-even projector
closure; local-current, EMT, and OAM consistency; and controlled replacement
of the analytic C3/C4 common-parent pilots.}
\end{quote}
The repository-specific prompt must use the actual H3 basis, current, TTN,
common-parent, Feshbach, and Wilson-handoff APIs.

\section{Risks, failure modes, and safeguards}
\label{sec:risks}

\begin{longtable}{L{0.22\textwidth}L{0.34\textwidth}L{0.36\textwidth}}
\toprule
Risk & Failure mode & Required safeguard\\
\midrule
\endhead
Incomplete color basis & One cluster tensor is called the five-parton color
space. & Compute the common generator nullspace and retain all three singlet
multiplicities.\\
Cluster statistics & Only quarks inside a chosen baryon cluster are
antisymmetrized. & Use exterior/occupation Fock space or an exact full
four-quark antisymmetrizer.\\
Nominal sea & Antiquarks are populated by a probability or copied curve. &
Require active positive-\(x\) antiquark slots and operator-derived overlaps.\\
Chiral overreach & A fitted contact kernel is called QCD chiral dynamics. &
Record it as a resolution-dependent effective operator with ablations,
flow, and holdouts.\\
Pion misidentification & A pseudoscalar pair is called an explicit pion. &
Require a pole/spectral and normalization matching statement.\\
PCAC double counting & \(2mP\) and \(f_\pi m_\pi^2\Phi\) are added as
separate divergence terms. & Use mutually exclusive WTI and pion-pole
representations joined by a matching cell.\\
Current inconsistency & Axial or pseudoscalar currents come from another
Hamiltonian. & Transform Hamiltonian and all operators together.\\
Tuned holdouts & \(\bar d-\bar u\), \(g_A\), or a TMD receives its own
coefficient. & Freeze holdouts and attach parameters only to shared physical
operators.\\
Energy-only TTN convergence & Low bond reproduces mass but loses sea/OAM or
PCAC. & Require observable-sensitive bond studies and protected symmetry
blocks.\\
Charge-conjugation error & One quark-minus-antiquark sign is used for every
operator. & Store and test the operator-specific \(C\)-signature.\\
Explicit/induced double counting & Pair sectors are added to the operator that
already represents them. & Enforce Feshbach equivalence plus remainder and
mutual exclusion.\\
Fake Wilson absorption & Numerical \(i\eps\) is interpreted as a physical
cut. & Require declared spectral support and exclude \(\eps\) from identity.\\
Premature phenomenology & Regulated H3 boundaries enter evolution or process
maps. & Maintain matching and factorization gates.\\
\bottomrule
\end{longtable}

\section{Recommended architecture and acceptance criteria}
\label{sec:acceptance}

\subsection{Recommended H3 foundation}

The adopted H3 foundation is
\begin{quote}
\emph{a fully antisymmetric, color-complete, four-sector light-front
Hamiltonian and fermionic symmetry-adapted tensor network, with canonical
pair creation, a controlled chiral pair operator, Hamiltonian-consistent
vector/axial/pseudoscalar currents, matched WTI and pion-pole PCAC
representations, direct positive-\(x\) antiquark GTMD overlaps, and a
provenance two-complex that makes explicit and induced sea descriptions
mutually exclusive except through a verified subtraction.}
\end{quote}

\subsection{Final formal acceptance criteria}

H3 is complete for its declared finite-basis scope only when:
\begin{enumerate}
 \item the state contains all four declared sectors with increasing basis
       support over at least three resolutions;
 \item all three five-parton color singlets are retained;
 \item exact four-quark antisymmetry holds before operator assembly;
 \item pair and chiral vertices obey all symmetry rules and Hermiticity;
 \item the H3 Hamiltonian is identical in dense and matrix-free action;
 \item sector-dependent parameters form a reproducible renormalization
       trajectory with exposed null directions;
 \item mass and vector charges close without fitting all holdouts;
 \item the H2 vector Ward benchmark remains passing;
 \item the nonsinglet axial WTI closes at the declared benchmark level;
 \item the pion-pole PCAC representation is matched to, rather than added to,
       the quark-level identity;
 \item every required axial/PCAC ablation gives a signed nonzero defect;
 \item exact, Krylov, and full-bond TTN states agree;
 \item variational TTN convergence is demonstrated in sea, OAM, axial, and
       PCAC observables;
 \item physical-state tracking is stable across the tower;
 \item probability, valence, charge, baryon, momentum, and canonical spin
       ledgers close;
 \item proton and neutron remain correlated members;
 \item positive-\(x\) \(\bar u\) and \(\bar d\) overlaps are evaluated from
       active antiquark slots;
 \item quark, antiquark, and gluon parents share one state identity;
 \item regulated TMD/GPD/PDF/current/EMT/OAM routes close;
 \item vector, axial, and tensor moments use their correct antiquark signs;
 \item explicit/induced sea equivalence includes transformed operators, norm
       kernel, and a visible remainder;
 \item the antiquark Wilson handoff retains complete identity and gives zero
       absorption without support;
 \item all assumption plans compile deterministically and incompatible plans
       fail;
 \item no production, nuclear, evolution, process, or inference object can
       consume H3 outputs;
 \item every released artifact includes the full assumption, convergence,
       provenance, matching-status, and unresolved-physics manifests.
\end{enumerate}

Completion of these criteria establishes the first common microscopic
quark--antiquark--gluon parent at finite resolution.  It does not by itself
establish continuum QCD, physical TMDs, or predictive deuteron phenomenology.

\appendix

\section{Compact mathematical and physical dictionary}

\begin{longtable}{L{0.30\textwidth}L{0.62\textwidth}}
\toprule
Physical concept & Formal object\\
\midrule
\endhead
Light-sea Fock support & \(\Hfiveu\oplus\Hfived\) with positive support and
flavor-resolved sector identity.\\
Five-parton color & Three-dimensional invariant subspace of
\(3^{\otimes4}\otimes\bar3\).\\
Four-quark statistics & Exterior power \(\bigwedge^4 h_q\) or exact
occupation-number representation.\\
Canonical pair creation & Typed \(g\leftrightarrow q\bar q\) Hamiltonian
block and adjoint.\\
Chiral pair mechanism & Resolution-dependent isovector scalar/pseudoscalar or
derivative effective operator.\\
Quark-level PCAC & Nonsinglet axial Ward identity
\(\partial\cdot A=2mP\).\\
Hadronic PCAC & Matched pion-pole plus non-pole representation of \(2mP\).\\
Microscopic antiquark & Positive-\(x\) active \(d^{\dagger}d\) contraction of
the bilocal quark operator.\\
Antiquark color transport & Anti-fundamental Wilson line
\(U_{\bar3}=U_3^*\).\\
Sea flavor asymmetry & Difference of direct \(\bar d\) and \(\bar u\)
projections from one correlated state.\\
Explicit/induced sea & Feshbach-equivalent descriptions with transformed
operators and a remainder.\\
Fermionic TTN & \(\mathbb Z_2\)-graded tensor network with swap signs and
symmetry blocks.\\
Common microscopic parent & Correlated vector \((W_q,W_{\bar q},W_g)\) with
one Hamiltonian/state/member identity.\\
Matching gate & Typed map from regulated H3 operators to a declared QCD
scheme; absent in H3.\\
Prediction & Withheld observable generated from shared parameters without
observable-specific retuning.\\
\bottomrule
\end{longtable}

\section{H3 assumption-bundle schema}

\begin{verbatim}
{
  "assumption_bundle_id": "...",
  "parent_plan": "H3-PLAN-A | H3-PLAN-B | H2-REFERENCE | INDUCED-SEA",
  "hamiltonian_route": "RSA_BLFQ_H3",
  "sectors": ["QQQ", "QQQG", "QQQUUBAR", "QQQDDBAR"],
  "resolution_tower": ["...", "...", "..."],
  "color_basis": {
    "qqq": 1,
    "qqqg": 2,
    "qqqq_qbar": 3,
    "recoupling_convention": "..."
  },
  "fermion_ordering": "...",
  "antisymmetry_id": "...",
  "canonical_pair_vertex": "enabled | disabled",
  "chiral_pair_vertex": "enabled | disabled",
  "pcac_representation": "QUARK_LEVEL_AXIAL_WTI | HADRONIC_PION_POLE",
  "confinement_branch": "refit | zero",
  "instantaneous_policy": "...",
  "zero_mode_policy": "...",
  "renormalization_conditions": ["..."],
  "holdout_manifest": "...",
  "ttn_topology": "...",
  "bond_policy": "...",
  "explicit_induced_sea_policy": "mutually_exclusive",
  "wilson_gate": "input_interface_only",
  "matching_gate": "closed",
  "nuclear_gate": "closed",
  "process_gate": "closed",
  "inference_gate": "closed",
  "source_hashes": {"...": "..."}
}
\end{verbatim}

\section{H3 microscopic state-bundle schema}

\begin{verbatim}
{
  "bundle_id": "...",
  "assumption_bundle_id": "...",
  "prediction_plan_id": "...",
  "hamiltonian_id": "...",
  "resolution_id": "...",
  "renormalization_trajectory_id": "...",
  "proton_neutron_member_id": "...",
  "state_tracking_id": "...",
  "sector_amplitudes": {
    "qqq": "...",
    "qqqg_rho": "...",
    "qqqg_lambda": "...",
    "qqqubar_color_1": "...",
    "qqqubar_color_rho": "...",
    "qqqubar_color_lambda": "...",
    "qqqdbar_color_1": "...",
    "qqqdbar_color_rho": "...",
    "qqqdbar_color_lambda": "..."
  },
  "exact_state_hash": "...",
  "krylov_state_hash": "...",
  "ttn_state_hash": "...",
  "fermionic_swap_ledger": "...",
  "bond_manifest_id": "...",
  "vector_current_bundle_id": "...",
  "axial_current_bundle_id": "...",
  "pseudoscalar_bundle_id": "...",
  "pcac_manifest_id": "...",
  "sea_flavor_ledger_id": "...",
  "common_parent_bundle_id": "...",
  "feshbach_comparison_id": "...",
  "antiquark_wilson_handoff_id": "...",
  "readiness": [
    "H3_LIGHT_SEA_BASIS_VALIDATED",
    "FINITE_BASIS_PCAC_BENCHMARKED",
    "POSITIVE_X_ANTIQUARK_OVERLAP_VALIDATED",
    "MICROSCOPIC_Q_QBAR_G_COMMON_PARENT_VALIDATED"
  ],
  "forbidden_readiness": [
    "PHYSICAL_TMD",
    "MATCHING_COMPLETE",
    "WILSON_READY",
    "NUCLEAR_MATCHING_READY",
    "INFERENCE_READY"
  ]
}
\end{verbatim}

\section{PCAC closure manifest}

\begin{longtable}{L{0.25\textwidth}L{0.67\textwidth}}
\toprule
Field & Required content\\
\midrule
\endhead
Identity & Hamiltonian, state, regulator, resolution, flavor generator, and
current source hashes.\\
Representation & \texttt{QUARK\_LEVEL\_AXIAL\_WTI} or
\texttt{HADRONIC\_PION\_POLE\_DECOMPOSITION}.\\
Kinematics & Incoming/outgoing momentum fibers, \(q^2\), current component,
and spinor convention.\\
One-body contribution & Active quark and antiquark current terms.\\
Pair contribution & Current attachment to canonical pair creation/annihilation
and five-parton sectors.\\
Chiral contribution & Current and pseudoscalar terms owned by the chiral
Hamiltonian kernel.\\
Instantaneous contribution & Constrained-field and inverse-derivative pieces.\\
Counterterms & Axial, pseudoscalar, wave-function, and current normalization
terms.\\
Pion-pole matching & Explicit or induced \(\Phi^a\), residue convention,
non-pole remainder, and matching map.\\
Residuals & Total and component defects with signs, tolerances, and numerical
error.\\
Ablations & Residual after removing each required term.\\
Holdouts & Second momentum point, \(g_A\), pion-pole, or GT diagnostic excluded
from calibration.\\
Status & Narrow benchmark passed and stronger claims forbidden.\\
\bottomrule
\end{longtable}

\section{Microscopic common-parent schema}

\begin{verbatim}
{
  "parent_bundle_id": "...",
  "state_bundle_id": "...",
  "member_id": "...",
  "resolution_id": "...",
  "recoil_map_id": "...",
  "parents": {
    "quark": ["u", "d"],
    "antiquark": ["ubar", "dbar"],
    "gluon": ["g"]
  },
  "operator_identities": ["..."],
  "active_slot_ancestry": ["..."],
  "color_outer_multiplicities": ["1", "rho", "lambda"],
  "helicity_oam_blocks": ["..."],
  "available_reductions": [
    "REGULATED_TMD",
    "REGULATED_GPD",
    "REGULATED_PDF",
    "VECTOR_CURRENT",
    "AXIAL_CURRENT",
    "PSEUDOSCALAR",
    "EMT",
    "WIGNER",
    "OAM"
  ],
  "matching_status": [
    "LINK_SHORTENING_REQUIRED",
    "UV_MATCHING_REQUIRED",
    "RAPIDITY_SOFT_MATCHING_REQUIRED",
    "NO_EVOLUTION_APPLIED",
    "NO_PROCESS_MAP_APPLIED"
  ],
  "unavailable_blocks": ["..."],
  "closure_report_id": "..."
}
\end{verbatim}

\section{Explicit/induced sea comparison schema}

\begin{verbatim}
{
  "comparison_id": "...",
  "explicit_state_bundle_id": "...",
  "retained_projector_id": "P",
  "eliminated_projector_id": "Q_pair",
  "energy_points": ["..."],
  "effective_hamiltonian_id": "...",
  "wave_operator_id": "...",
  "norm_kernel_id": "...",
  "transformed_operator_ids": {
    "vector": "...",
    "axial": "...",
    "pseudoscalar": "...",
    "antiquark_bilocal": "...",
    "gtmd": "..."
  },
  "equivalence_residuals": {"...": 0.0},
  "remainder_norms": {"...": "nonzero"},
  "conditioning": {"...": "..."},
  "provenance_two_cell": "...",
  "composition_rule": "EXPLICIT_XOR_INDUCED"
}
\end{verbatim}

\section{Antiquark Wilson-handoff schema}

\begin{verbatim}
{
  "handoff_id": "...",
  "state_bundle_id": "...",
  "common_parent_bundle_id": "...",
  "active_species": "ubar | dbar",
  "active_helicity": "...",
  "pair_sector_ancestry": "...",
  "five_parton_color_multiplicity": "1 | rho | lambda",
  "anti_fundamental_path_id": "...",
  "charge_conjugation_adapter_id": "...",
  "oam_blocks": [0, -1, 1, -2, 2],
  "hamiltonian_coupling_id": "...",
  "cut_support_status": "ABSENT | FINITE_VOLUME | CONTINUUM",
  "phase_budget_id": "...",
  "soft_overlap_status": "...",
  "wilson_order": 1,
  "readiness": "MICROSCOPIC_ANTIQUARK_WILSON_INPUT_INTERFACE_VALIDATED"
}
\end{verbatim}

\section{Negative-test catalogue}

The implementation shall assign stable IDs to at least the following classes.
They are grouped here to make the failure surface auditable rather than to
limit the actual test count.

\begin{longtable}{L{0.18\textwidth}L{0.72\textwidth}}
\toprule
Group & Required injected failures\\
\midrule
\endhead
Color & Missing each of the three singlets; nonorthogonal basis; wrong
anti-fundamental generator; cluster tensor promoted to complete basis; lost
outer multiplicity; nonunitary recoupling.\\
Statistics & Wrong exchange sign; post-assembly antisymmetry; Pauli-forbidden
state accepted; cluster-only antisymmetry; missing fermionic swap sign.\\
Kinematics & Nonpositive support; momentum nonclosure; wrong charge, baryon,
isospin, \(\Jz\), center-of-mass, recoil, or zero-skewness identity.\\
Hamiltonian & Missing pair adjoint; wrong flavor/color/helicity selection;
unsupported block inserted; sector counterterm misassigned; chiral kernel
called universal QCD; H1 color-spin added to explicit H2/H3 dynamics.\\
Renormalization & Bare sector mass called physical; resolution scale reused as
TMD scale; holdout leakage; null direction hidden; sea asymmetry fitted with a
new coefficient.\\
Currents & Vector current from wrong Hamiltonian; pair, chiral, instantaneous,
pseudoscalar, pion-pole, or counterterm piece omitted; convention mismatch;
one-point PCAC advertised as a continuum proof; singlet anomaly ignored.\\
Solvers and TTN & Exact/Krylov mismatch; matrix-free mismatch; full-bond
mismatch; variational energy below exact; nonmonotone nested optimum; required
sea/color block removed; low-rank loss unreported; state branch swapped.\\
Ledgers & Probability, valence, charge, baryon, momentum, spin, or
proton/neutron member closure broken; global \(\bar u=\bar d\) imposed.\\
Antiquark parent & Negative-\(x\) copy; wrong active species; duplicate active
multiplicity; wrong charge-conjugation sign; off-diagonal overlap without
source; local recoil reimplementation.\\
Reductions & Quark number convention used for gluon; wrong quark/antiquark sign
in vector, axial, or tensor moment; WTI/PCAC routes added; false matching
completion; physical PDF/GTMD/TMD status asserted.\\
Feshbach and provenance & Operator not transformed; norm kernel omitted;
remainder hidden; singular resolvent undetected; explicit plus induced sea;
missing two-cell; unresolved cycle ignored; explicit pair plus independent
pion-cloud correction.\\
Wilson and gates & Antiquark flavor/color/OAM lost; numerical epsilon made
physical; absorption without support; C5/C6 types duplicated; evolution,
process, nuclear, inference, or production promotion.\\
Regression & Production registry, parent artifact, provenance, composition,
prior microscopic oracle, pinned Wilson manifest, or normative source changed.\\
\bottomrule
\end{longtable}

\small
\begin{thebibliography}{99}
\raggedright

\bibitem{Dirac1949}
P.~A.~M.~Dirac,
``Forms of relativistic dynamics,''
Rev. Mod. Phys. \textbf{21}, 392 (1949).

\bibitem{BrodskyPauliPinsky1998}
S.~J.~Brodsky, H.-C.~Pauli, and S.~S.~Pinsky,
``Quantum chromodynamics and other field theories on the light cone,''
Phys. Rept. \textbf{301}, 299 (1998),
arXiv:hep-ph/9705477.

\bibitem{Vary2010}
J.~P.~Vary et al.,
``Hamiltonian light-front field theory in a basis function approach,''
Phys. Rev. C \textbf{81}, 035205 (2010),
arXiv:0905.1411.

\bibitem{Xu2025}
S.~Xu, Y.~Liu, C.~Mondal, J.~Lan, X.~Zhao, Y.~Li, and J.~P.~Vary,
``Towards a first principles light-front Hamiltonian for the nucleon,''
Phys. Lett. B \textbf{867}, 139599 (2025),
arXiv:2408.11298.

\bibitem{Mondal2025}
C.~Mondal, S.~Xu, Y.~Liu, J.~Lan, Y.~Li, X.~Zhao, and J.~P.~Vary,
``Basis light-front quantization: Advancing a first principles approach for
the nucleon,''
arXiv:2503.21372 (2025).

\bibitem{Xu2021}
S.~Xu, C.~Mondal, J.~Lan, X.~Zhao, Y.~Li, and J.~P.~Vary,
``Nucleon structure from basis light-front quantization,''
Phys. Rev. D \textbf{104}, 094036 (2021),
arXiv:2108.03909.

\bibitem{Du2020}
W.~Du, Y.~Li, X.~Zhao, G.~A.~Miller, and J.~P.~Vary,
``Basis light-front quantization for a chiral nucleon-pion Lagrangian,''
Phys. Rev. C \textbf{101}, 035202 (2020),
arXiv:1911.10762.

\bibitem{Duan2024}
Y.~Duan, S.~Xu, S.~Cheng, X.~Zhao, Y.~Li, and J.~P.~Vary,
``Flavor asymmetry from the non-perturbative nucleon sea,''
arXiv:2404.07755 (2024).

\bibitem{Cao2026}
X.~Cao, S.~Cheng, Y.~Duan, Y.~Li, S.~Xu, and X.~Zhao,
``Non-perturbative flavor asymmetry in the nucleon and deuteron: The
light-front Hamiltonian effective field theory approach,''
arXiv:2601.13567 (2026).

\bibitem{KoflerPasquini2017}
S.~Kofler and B.~Pasquini,
``Collinear parton distributions and the structure of the nucleon sea in a
light-front meson-cloud model,''
Phys. Rev. D \textbf{95}, 094015 (2017),
arXiv:1701.07839.

\bibitem{Choudhary2023}
P.~Choudhary, D.~Chakrabarti, and C.~Mondal,
``Flavor asymmetry of light sea quarks in proton: A light-front spectator
model,''
arXiv:2312.01484 (2023).

\bibitem{NJL1961}
Y.~Nambu and G.~Jona-Lasinio,
``Dynamical model of elementary particles based on an analogy with
superconductivity. I,''
Phys. Rev. \textbf{122}, 345 (1961).

\bibitem{WuZhang2004}
M.-H.~Wu and W.-M.~Zhang,
``Chiral symmetry in light-front QCD,''
JHEP \textbf{04}, 045 (2004),
arXiv:hep-ph/0310095.

\bibitem{Nambu1960}
Y.~Nambu,
``Axial vector current conservation in weak interactions,''
Phys. Rev. Lett. \textbf{4}, 380 (1960).

\bibitem{GoldbergerTreiman1958}
M.~L.~Goldberger and S.~B.~Treiman,
``Decay of the pi meson,''
Phys. Rev. \textbf{110}, 1178 (1958),
and
``Form factors in beta decay and muon capture,''
Phys. Rev. \textbf{111}, 354 (1958).

\bibitem{Chen2021}
C.~Chen, C.~S.~Fischer, C.~D.~Roberts, and J.~Segovia,
``Nucleon axial-vector and pseudoscalar form factors, and PCAC relations,''
Phys. Rev. D \textbf{105}, 094022 (2022),
arXiv:2103.02054.

\bibitem{Alexandrou2024}
C.~Alexandrou et al.,
``Nucleon axial and pseudoscalar form factors using twisted-mass fermion
ensembles at the physical point,''
Phys. Rev. D \textbf{109}, 034503 (2024),
arXiv:2309.05774.

\bibitem{Bali2019}
G.~S.~Bali et al.,
``Solving the PCAC puzzle for nucleon axial and pseudoscalar form factors,''
Phys. Lett. B \textbf{789}, 666 (2019),
arXiv:1810.05569.

\bibitem{BrodskyDiehlHwang2001}
S.~J.~Brodsky, M.~Diehl, and D.-S.~Hwang,
``Light-cone wavefunction representation of deeply virtual Compton
scattering,''
Nucl. Phys. B \textbf{596}, 99 (2001),
arXiv:hep-ph/0009254.

\bibitem{Meissner2009}
S.~Mei{\ss}ner, A.~Metz, and M.~Schlegel,
``Generalized parton correlation functions for a spin-1/2 hadron,''
JHEP \textbf{08}, 056 (2009),
arXiv:0906.5323.

\bibitem{Lorce2011}
C.~Lorc\'e and B.~Pasquini,
``Quark Wigner distributions and orbital angular momentum,''
Phys. Rev. D \textbf{84}, 014015 (2011),
arXiv:1106.0139.

\bibitem{Feshbach1958}
H.~Feshbach,
``Unified theory of nuclear reactions,''
Ann. Phys. \textbf{5}, 357 (1958),
and Ann. Phys. \textbf{19}, 287 (1962).

\bibitem{Karmanov2008}
V.~A.~Karmanov, J.-F.~Mathiot, and A.~V.~Smirnov,
``Systematic renormalization scheme in light-front dynamics with Fock-space
truncation,''
Phys. Rev. D \textbf{77}, 085028 (2008),
arXiv:0801.4507.

\bibitem{Marinho2008}
J.~A.~O.~Marinho, T.~Frederico, E.~Pace, G.~Salm\`e, and P.~U.~Sauer,
``Light-front Ward--Takahashi identity for two-fermion systems,''
Phys. Rev. D \textbf{77}, 116010 (2008),
arXiv:0805.0707.

\bibitem{Singh2010}
S.~Singh, R.~N.~C.~Pfeifer, and G.~Vidal,
``Tensor network decompositions in the presence of a global symmetry,''
Phys. Rev. A \textbf{82}, 050301(R) (2010),
arXiv:0907.2994.

\bibitem{Corboz2010}
P.~Corboz, G.~Evenbly, F.~Verstraete, and G.~Vidal,
``Simulation of interacting fermions with entanglement renormalization,''
Phys. Rev. A \textbf{81}, 010303(R) (2010),
arXiv:0904.4151.

\bibitem{Gardi2013}
E.~Gardi, J.~M.~Smillie, and C.~D.~White,
``The non-Abelian exponentiation theorem for multiple Wilson lines,''
JHEP \textbf{06}, 088 (2013),
arXiv:1304.7040.

\bibitem{Collins2011}
J.~Collins,
\emph{Foundations of Perturbative QCD},
Cambridge University Press (2011).

\end{thebibliography}

\end{document}
